% NOTE: uncomment if using arXiv format, comment block below otherwise
% \documentclass[11pt]{article}


% % NOTE: uncomment if using Biometrika format, comment block below otherwise
% %% FIX INCOMPATIBILITY OF BIOMETRIKA CLASS WITH LATEX KERNEL
% %% https://tex.stackexchange.com/questions/575327/extra-endgroup-when-using-the-biometrika-template
% %% save the original kernel definitions
\let\latexarabic\arabic
\let\latexdocument\document
\let\latexenddocument\enddocument
%%% fix for bad usage of ntheorem
\RequirePackage[thmmarks]{ntheorem}
\makeatletter
\renewtheoremstyle{plain}
  {\item[\hskip\labelsep \theorem@headerfont ##1\ \textup{##2}\theorem@separator]}
  {\item[\hskip\labelsep \theorem@headerfont ##1\ \textup{##2}\ (##3)\theorem@separator]}
\makeatother
\documentclass[lineno]{biometrika}


% load various specifications
% standard packages and options
\usepackage{graphicx}
\usepackage{float}
\usepackage{lmodern}
\usepackage{booktabs}
\usepackage{amsmath,amsfonts,amssymb,bbm,bm,dsfont,mathtools,relsize,exscale}
\usepackage{tikz}
\usepackage[dvipsnames]{xcolor}
\usepackage[OT1]{fontenc}
\usepackage[utf8]{inputenc}
\usepackage[linesnumbered,ruled,vlined]{algorithm2e}
\usepackage[inline,shortlabels]{enumitem}
\bibliographystyle{biometrika}
\graphicspath{ {figs/} }

% for cross-linking documents
\usepackage{refcount,nameref,zref-xr,zref-user} % nameref before zref-xr
\zxrsetup{toltxlabel} % allows use of the \ref command as usual
\usepackage{url}

% https://tex.stackexchange.com/questions/339408/how-to-color-cross-reference-with-oxford-bioinfomatics-latex-template
\let\href\undefined
\usepackage[colorlinks,citecolor=blue,urlcolor=blue]{hyperref}

% journal-specific packages
\makeatletter%
\@ifclassloaded{biometrika}{%
  %% FIX INCOMPATIBILITY OF BIOMETRIKA CLASS WITH LATEX KERNEL
  %%% recover the original definitions
  \let\document\latexdocument
  \let\enddocument\latexenddocument
  \AtEndDocument{\printhistory}
  \let\arabic\latexarabic
  \def\rm{}

  %%% User-defined macros should be placed here, but keep them to a minimum.
  \def\Bka{{\it Biometrika}}
  \def\AIC{\textsc{aic}}
  \def\T{{ \mathrm{\scriptscriptstyle T} }}
  \def\v{{\varepsilon}}

  %\addtolength\topmargin{35pt}
  \DeclareMathOperator{\Thetabb}{\mathcal{C}}

  % authorship details
  \jname{Biometrika}
  %% The year, volume, and number are determined on publication
  %\jyear{2022}
  %\jvol{X}
  %\jnum{X}

  %% The \doi{...} and \accessdate commands are used by the production team
  %\doi{10.1093/biomet/asm023}
  %\accessdate{Advance Access publication on 30 August 2020}

  %% These dates are usually set by the production team
  %\received{15 May 2022}
  %\revised{15 November 2022}

  %% The left and right page headers are defined here:
  \markboth{Testa et al.}{Efficient inference for conditional longitudinal modified treatment policies}

  \newcommand{\authorlist}{
      %% Here are the title, author names and addresses
      \author{Christian Testa} 
      \affil{Department of Biostatistics,\\
        Harvard T.H. Chan School of Public Health,\\
        665 Huntington Avenue, Boston, MA 02115
      \email{ctesta@hsph.harvard.edu}}

      \author{Anybody who wants to join}
      \affil{Department ...,\\
        Institute,\\
        665 Huntington Avenue, Boston, MA 02115
      \email{}}

      \author{Christopher D. Golden}
      \affil{Department of Nutrition, \\
        Department of Environmental Health,
        Harvard T.H. Chan School of Public Health,\\
        665 Huntington Avenue, Boston, MA 02115\\
        \email{golden@hsph.harvard.edu}
      \email{}}

      \author{Nima S. ~Hejazi}
      \affil{Department of Biostatistics, \\
        Harvard T.H. Chan School of Public Health,\\
        665 Huntington Avenue, Boston, MA 02115
      \email{nhejazi@hsph.harvard.edu}}

  }
}{%
  %for plain article format only
  \usepackage{arxiv}
  \usepackage{standalone}
  \usepackage{multirow}
  \usepackage{setspace}
  \usepackage{amsthm}
  \usepackage[round]{natbib}

  % custom math definitions and macros that conflict with Biometrika
  \newtheorem{theorem}{Theorem}
  \newtheorem{lemma}{Lemma}
  \newtheorem{remark}{Remark}
  \newtheorem{definition}{Definition}
  \newtheorem{proposition}{Proposition}
  {\theoremstyle{definition}\newtheorem{assumption}{}}
  {\theoremstyle{definition}\newtheorem{condition}{}}
  {\theoremstyle{definition}\newtheorem{example}{Example}[section]}

  % authorship details
  \newcommand{\authorlist}{
    Christian Testa \\
    Department of Biostatistics,\\
    Harvard T.H. Chan School of Public Health,\\
    \texttt{ctesta@hsph.harvard.edu}\\
    \And
    Anybody who wants to join \\
    Department ...,\\
    Institute ...,\\
    \texttt{email}\\
    \And
    Christopher D.~Golden \\
    Department of Nutrition, Department of Environmental Health,\\
    Harvard T.H. Chan School of Public Health,\\
    \texttt{golden@hsph.harvard.edu}\\
    \And
    Nima S.~Hejazi \\
    Department of Biostatistics,\\
    Harvard T.H. Chan School of Public Health,\\
    \texttt{nhejazi@hsph.harvard.edu}\\
  }
}
\makeatother

% paper title
\newcommand{\titlepaper}{Efficient inference for conditional longitudinal modified treatment policies}

% counters and definitions
\AtEndDocument{\refstepcounter{theorem}\label{finalthm}}
\AtEndDocument{\refstepcounter{equation}\label{finaleq}}
\AtEndDocument{\refstepcounter{lemma}\label{finallemma}}
\renewcommand\theassumption{{A}\arabic{assumption}}
\renewcommand\thecondition{{C}\arabic{condition}}
\renewcommand\theexample{\arabic{example}}

% math macros
\renewcommand{\d}{\mathbbm{d}}
\newcommand{\D}{\mathcal{D}}
\newcommand{\E}{\mathbb{E}}
\newcommand{\F}{\mathcal{F}}
\newcommand{\M}{\mathcal{M}}
\newcommand{\R}{\mathbb{R}}
\newcommand{\X}{\mathcal{X}}
\newcommand{\logit}{\text{logit}}
\renewcommand{\P}{\mathbb{P}}
\newcommand{\prob}{\mathbb{P}}
\newcommand{\1}{\mathbbm{1}}
\newcommand{\indep}{\mbox{$\perp\!\!\!\perp$}}
\DeclareMathOperator*{\argmin}{\arg\!\min}
\DeclareMathOperator*{\argmax}{\arg\!\max}

% editors
\usepackage[dvipsnames]{xcolor}
\newcommand{\nh}[1]{\textcolor{Cyan}{*** Nima: #1}}
\newcommand{\ct}[1]{\textcolor{ForestGreen}{*** Christian: #1}}

\zexternaldocument*[supp:]{sm}


\makeatletter
% FOR BIOMETRIKA
\@ifclassloaded{biometrika}{%
  % set list of authors
  \authorlist
}{%
%% FOR ARXIV
  \author{\authorlist}
  \date{\today}
}
\makeatother
\title{\titlepaper}

\begin{document}
\maketitle


\begin{abstract}
Longitudinal modified treatment policies (LMTPs) are an increasingly common
approach for the estimation and inference of causal effects corresponding to changes in longitudinal exposures or treatment. In many problems, scientific interest lies not only in 
overall effects but also subgroup-specific or conditional effects. However, little work has
been done to address efficiently estimating the effects of LMTPs on many
subgroups simultaneously. We present a novel targeted learning approach for the 
efficient estimation of subgroup-specific LMTP effects, where subgroups are defined by characteristics of observations at baseline,
which provides multiply robust and semiparametric efficient estimators. 
Our procedure for estimating subgroup effects of LMTPs allows for the construction of
simultaneous confidence intervals across subgroups. For handling 
continuous covariates defined at baseline, we develop a DR-learner inspired approach for 
estimating conditional LMTPs consisting of forming an efficient influence function based
pseudo-outcome regressed on baseline variables.  
As a case study, we apply 
our method to the question of how deforestation in Madagascar during the 2010-2022 timeframe
may have affected malaria incidence rates at 2,773 clinics where effect heterogeneity was expected
by ecological zone.

Keywords: Causal inference, Longitudinal Modified Treatment Policies, Treatment Effect
Heterogeneity, Efficient Inference, Subgroup Effects
\end{abstract}


%%%%%%%%%%%%%%%%%%%%%%%%%%%%%%%%%%%%%%%%%%%%%%%%%%%%%%%%%%%%%%%%%%%%%%%%%%%%%%%
\section{Introduction}\label{intro}

Longitudinal studies in medicine, epidemiology, and the environmental sciences 
often seek to evaluate interventions on exposures or treatments that evolve over time. 
In such settings, treatment decisions may depend on a subject's previous treatment and
covariate history, and time-varying covariates may themselves be affected by earlier
treatment.  These features rule out standard regression adjustment and motivate 
the use of causal estimands defined through longitudinal interventions. 

Longitudinal modified treatment policies (LMTPs) provide one useful
class of such interventions. Rather than setting treatment to one fixed value, 
an LMTP modifies the treatment that a subject would otherwise have received according
to a user-specified policy map. This makes LMTPs attractive in settings where
deterministic static interventions are implausible or scientifically less relevant. For example,
in studies of environmental exposures, one may be interested in shifting an exposure 
downward for only units above a threshold; and in health studies, modifying a
dose or treatment intensity relative to the value actually received. LMTPs therefore
provide a flexible way to encode realistic longitudinal interventions while retaining
identification under standard sequential causal assumptions. 

Most methodological work on LMTPs has focused on marginal causal parameters,
such as the population mean outcome under a single policy or the 
contrast between two policies. However, scientific questions about policy effects are often
inherently conditional. Investigators may ask whether an intervention has 
larger effects in certain ecological zones, clinical subgroups, baseline risk strata,
or other baseline-defined populations. More generally, one may wish
to estimate a conditional effect curve describing how the causal effect of a longitudinal 
policy varies with baseline covariates.  Restricting attention to only marginal 
LMTP effects can obscure such heterogeneity. 

This paper studies estimation and inference for conditional effects of LMTPs. 
Let $V$ denote baseline covariates with respect to which 
effect heterogeneity is of interest. For an LMTP $\mathbbm d$, we define the conditional 
LMTP mean 
\[ 
\theta_{\mathbbm d}(v) =   \mathbb{E}\{Y(\mathbbm d) \mid V=v\},
\]
and, for two policies \(\mathbbm d_1\) and \(\mathbbm d_0\), the conditional LMTP
contrast
\[
  \tau(v)
  =
  \theta_{\mathbbm d_1}(v) - \theta_{\mathbbm d_0}(v).
\]
When $V$ is discrete, this target corresponds to a vector of subgroup-specific
policy effects. When $V$ is continuous or high-dimensional, it corresponds to a
conditional effect function. 

We will first present a formal identification
argument for conditional LMTP means and contrasts using the observed-data
likelihood factorization, the post-policy law, and the longitudinal
\(g\)-formula. This yields inverse-probability and iterated-regression
representations of the conditional LMTP effect. Second, for finite baseline
subgroups, we develop a targeted learning estimator for the vector of subgroup
LMTP means or contrasts. This estimator solves the efficient influence-function
estimating equations for all subgroups jointly and supports Wald-type and
simultaneous inference. Third, for general conditioning covariates \(V\), we
develop a DR-Learner for conditional LMTP effects. The DR-Learner estimates
uncentered influence-function pseudo-outcomes for marginal LMTP parameters and
then regresses these pseudo-outcomes on \(V\). We establish conditions under
which this estimator is asymptotically equivalent to the benchmark second-stage
regression that would be obtained if the nuisance functions were known.

The remainder of the paper is organized as follows. Section~\ref{sec:setup}
defines the observed data structure, LMTPs, conditional target parameters,
identification assumptions, identification results and derives the
pseudo-outcome representation used by the DR-Learner.
Section~\ref{sec:estimation} presents the finite-subgroup targeted minimum
loss-based estimator and the conditional LMTP DR-Learner.
Section~\ref{sec:theory} gives the large-sample properties of the proposed
estimators. Section~\ref{sec:simulations} presents simulation studies, and
Section~\ref{sec:application} applies the methods to estimate heterogeneous
effects of deforestation on malaria incidence in Madagascar.
Section~\ref{sec:discussion} offers discussion of potential applications for
conditional LMTP estimation and future directions for conditional effect
estimation in longitudinal causal inference.

\section{Setup and Target Parameters}
\label{sec:setup}
\subsection{Observed data and likelihood factorization}
\label{subsec:observed-data}

Let
\[
  O = (L_0, A_0, L_1, A_1, \ldots, L_\tau, A_\tau, Y)
\]
denote the observed data for a randomly sampled unit. Here \(L_0\) denotes
baseline covariates, \(A_t\) denotes the treatment or exposure at time
\(t=0,\ldots,\tau\), and \(L_{t+1}\) denotes covariates measured after treatment
\(A_t\), with \(L_{\tau+1} \equiv Y\). Let
\[
  \overline L_t = (L_0,\ldots,L_t),
  \qquad
  \overline A_t = (A_0,\ldots,A_t),
\]
and define the treatment history available immediately before treatment time
\(t\) by
\[
  H_t = (\overline L_t,\overline A_{t-1}),
\]
with the convention that \(H_0=L_0\). We assume throughout that
\(O_1,\ldots,O_n\) are independent and identically distributed according to an
unknown law \(P\).

Assume that \(P\) is dominated by a product measure and has density
\[
  p(o)
  =
  q_0(l_0)
  \prod_{t=0}^{\tau}
  g_t(a_t \mid h_t)
  q_{t+1}(l_{t+1} \mid h_t,a_t),
\]
where
\[
  g_t(a_t\mid h_t)
  =
  p(A_t=a_t\mid H_t=h_t)
\]
denotes the conditional density or probability mass function of the treatment at
time \(t\), and
\[
  q_{t+1}(l_{t+1}\mid h_t,a_t)
  =
  p(L_{t+1}=l_{t+1}\mid H_t=h_t,A_t=a_t)
\]
denotes the conditional law of the next covariate, or of the final outcome when
\(t=\tau\). We write
\[
  q = (q_0,q_1,\ldots,q_{\tau+1}),
  \qquad
  g = (g_0,\ldots,g_\tau),
\]
so that \(p=qg\). This factorization separates the conditional treatment process
\(g\) from the covariate and outcome transition laws \(q\). The causal parameters
considered below are identified by replacing the observed treatment process
\(g\) with a post-policy treatment process while leaving \(q\) unchanged.

\subsection{Longitudinal modified treatment policies}
\label{subsec:lmtps}

A longitudinal modified treatment policy is a sequence
\[
  \mathbbm d = (d_0,\ldots,d_\tau)
\]
of known maps. At time \(t\), the map \(d_t\) takes an observed treatment value
and history \((a_t,h_t)\) and returns the treatment value that would be assigned
under the policy,
\[
  a_t^{\mathbbm d} = d_t(a_t,h_t).
\]
The corresponding random variable is
\[
  A_t^{\mathbbm d} = d_t(A_t,H_t).
\]
For example, with a continuous exposure, \(d_t\) might shift exposure downward by
a fixed amount among units above a threshold. With a discrete exposure, \(d_t\)
might map some observed treatment levels to modified levels while leaving other
levels unchanged.

In the longitudinal setting, the policy is evaluated recursively. Earlier
modified treatments affect later histories, and later modified treatments are
therefore functions of histories that include previously modified treatment
values. Let
\[
  H_t^{\mathbbm d}
  =
  (\overline L_t, A_0^{\mathbbm d},\ldots,A_{t-1}^{\mathbbm d})
\]
denote the history through time \(t\) under the policy. Then the policy-modified
treatment process is generated recursively by
\[
  A_0^{\mathbbm d}
  =
  d_0(A_0,L_0),
\]
and, for \(t=1,\ldots,\tau\),
\[
  A_t^{\mathbbm d}
  =
  d_t\{A_t(\overline A_{t-1}^{\mathbbm d}),H_t^{\mathbbm d}\},
\]
where \(A_t(\overline A_{t-1}^{\mathbbm d})\) denotes the treatment value that
would have been observed at time \(t\) had treatment through time \(t-1\) been
modified according to \(\mathbbm d\).

Let
\[
  g_t^{\mathbbm d}(a_t\mid h_t)
\]
denote the conditional law of \(A_t^{\mathbbm d}\) given \(H_t=h_t\). For
discrete treatments, \(g_t^{\mathbbm d}\) is the pushforward of the observed
treatment law \(g_t\) under the map \(d_t\). For continuous treatments, we assume
that \(d_t(\cdot,h_t)\) is piecewise smooth and invertible on a finite partition
of the treatment space, so that \(g_t^{\mathbbm d}\) is well-defined by the
change-of-variables formula.

The post-policy distribution \(P^{\mathbbm d}\) is the law with density
\[
  p^{\mathbbm d}(o)
  =
  q_0(l_0)
  \prod_{t=0}^{\tau}
  g_t^{\mathbbm d}(a_t\mid h_t)
  q_{t+1}(l_{t+1}\mid h_t,a_t).
\]
Thus the intervention changes the conditional treatment law at each time point
but leaves the conditional covariate and outcome transition laws unchanged. This
is the usual longitudinal \(g\)-formula representation specialized to modified
treatment policies.

\subsection{Conditional LMTP means and contrasts}
\label{subsec:targets}

Let \(V\subseteq L_0\) denote baseline covariates with respect to which effect
heterogeneity is of interest. The covariates \(V\) may be a discrete subgroup
indicator, a vector of baseline features, or a lower-dimensional summary of
baseline risk. Because \(V\) is measured before treatment, its distribution is not
changed by the intervention.

For a single LMTP \(\mathbbm d\), define the marginal LMTP mean
\[
  \psi_{\mathbbm d}
  =
  \mathbb{E}_{P^{\mathbbm d}}(Y)
  =
  \mathbb{E}\{Y(\mathbbm d)\}.
\]
The corresponding conditional LMTP mean is
\[
  \theta_{\mathbbm d}(v)
  =
  \mathbb{E}_{P^{\mathbbm d}}(Y\mid V=v)
  =
  \mathbb{E}\{Y(\mathbbm d)\mid V=v\}.
\]
For two policies \(\mathbbm d_1\) and \(\mathbbm d_0\), define the conditional LMTP contrast
\[
  \tau(v)
  =
  \theta_{\mathbbm d_1}(v)
  -
  \theta_{\mathbbm d_0}(v) = 
  \E\{Y^{\mathbbm d_1} - Y^{\mathbbm d_0} \mid V = v\}.
\]

When \(V\) is discrete with support \(\mathcal V\), the target is the finite
collection
\[
  \{\theta_{\mathbbm d}(v):v\in\mathcal V\}
\]
or, for a contrast of two policies,
\[
  \{\tau(v):v\in\mathcal V\}.
\]
In this case, the target is a finite-dimensional parameter, and standard
semiparametric efficiency theory can be used to construct asymptotically linear
estimators and Wald-type confidence intervals.

When \(V\) is continuous, the target \(v\mapsto \theta_{\mathbbm d}(v)\) or
\(v\mapsto \tau(v)\) is an infinite-dimensional conditional mean function. In
this case, the conditional effect itself is not generally a pathwise
differentiable parameter at a point \(v\) without additional smoothing or
projection. We therefore distinguish between finite-dimensional subgroup
inference and flexible estimation of conditional effect functions.

\subsection{Identification assumptions}
\label{subsec:identification-assumptions}

Identification of the conditional LMTP mean relies on a longitudinal causal
model, a support condition for the modified treatment policy, and regularity
conditions ensuring that the post-policy treatment density is well-defined.

We first introduce a nonparametric structural equation model to state the
sequential randomization condition precisely. Suppose the observed data are
generated by structural equations
\[
  L_0 = f_{L_0}(U_{L_0}),
\]
and, for \(t=0,\ldots,\tau\),
\[
  A_t = f_{A_t}(H_t,U_{A_t}),
  \qquad
  L_{t+1} = f_{L_{t+1}}(H_t,A_t,U_{L_{t+1}}),
\]
where \(L_{\tau+1}=Y\). For \(t=0,\ldots,\tau\), define
\[
  \underline U_{L,t+1}
  =
  (U_{L_{t+1}},\ldots,U_{L_{\tau+1}}),
  \qquad
  \underline U_{A,t+1}
  =
  (U_{A_{t+1}},\ldots,U_{A_\tau}).
\]
The exogenous variables encode sources of variation not included in the observed
history.

\begin{assumption}[Consistency]
\label{assump:consistency}
The observed data are generated by the structural equations above. For any
treatment history \(\overline a_\tau\), the counterfactual outcome
\(Y(\overline a_\tau)\) is obtained by replacing the treatment equations with
\(A_t=a_t\), \(t=0,\ldots,\tau\), while leaving the covariate and outcome
equations unchanged. In particular, for the observed treatment history,
\[
  Y = Y(\overline A_\tau).
\]
% Likewise, each observed post-treatment covariate \(L_{t+1}\) agrees with its
% counterfactual value under the observed treatment history through time \(t\).
\end{assumption}

\begin{assumption}[Strong sequential randomization]
\label{assump:seq-randomization}
For each \(t=0,\ldots,\tau\),
\[
  U_{A_t}
  \perp\!\!\!\perp
  \left(
    \underline U_{L,t+1},
    \underline U_{A,t+1}
  \right)
  \mid H_t.
\]
Equivalently, conditional on the observed history \(H_t\), the unobserved
determinants of treatment at time \(t\) are independent of the unobserved
determinants of future covariates, future treatments, and the final outcome.
\end{assumption}

\begin{remark}
Assumption~\ref{assump:seq-randomization} is the NPSEM version of no unmeasured
time-varying confounding. It implies the more familiar sequential
exchangeability statement that, conditional on \(H_t\), the treatment actually
received at time \(t\) is independent of future counterfactual variables under
the treatment regimes considered in the analysis.
\end{remark}

\begin{assumption}[LMTP support condition]
\label{assump:lmtp-support}
For each \(t=0,\ldots,\tau\), if
\[
  (a_t,h_t) \in \operatorname{supp}(A_t,H_t),
\]
then
\[
  \{d_t(a_t,h_t),h_t\}
  \in
  \operatorname{supp}(A_t,H_t).
\]
Equivalently, for \(P\)-almost every history \(h_t\), the modified treatment
value \(d_t(a_t,h_t)\) lies in the conditional support of \(A_t\mid H_t=h_t\)
whenever \(a_t\) lies in that support.
\end{assumption}

\begin{remark}
Assumption~\ref{assump:lmtp-support} is the LMTP analogue of positivity. It is
weaker and more local than the usual static-treatment positivity requirement.
It does not require every treatment level to be possible at every history.
Instead, it requires only that the treatment value assigned by the modified
policy be observable, in principle, among units with the same history.
\end{remark}

\begin{assumption}[Regularity of the policy map]
\label{assump:policy-regularity}
For each \(t=0,\ldots,\tau\), the map \(d_t(\cdot,h_t)\) is known. If \(A_t\) is
continuous, then \(d_t(\cdot,h_t)\) is piecewise smooth and invertible on a
finite partition of the treatment space for \(P\)-almost every \(h_t\). On each
piece, the inverse map \(d_t^{-1}(\cdot,h_t)\) exists and has derivative bounded
away from zero and infinity on the relevant support. If \(A_t\) is discrete, the
pushforward law induced by \(d_t\) is well-defined.
\end{assumption}

Under Assumption~\ref{assump:policy-regularity}, the post-policy treatment
density is defined as the law of \(A_t^{\mathbbm d}=d_t(A_t,H_t)\) conditional on
\(H_t\). For continuous treatments, this density is obtained by the
change-of-variables formula. In the one-to-one case,
\[
  g_t^{\mathbbm d}(a_t\mid h_t)
  =
  g_t\{d_t^{-1}(a_t,h_t)\mid h_t\}
  \left|
    \frac{\partial}{\partial a_t}
    d_t^{-1}(a_t,h_t)
  \right|.
\]
For piecewise invertible maps, the density is the sum of the corresponding
piecewise inverse contributions. For discrete treatments,
\[
  g_t^{\mathbbm d}(a_t\mid h_t)
  =
  \sum_{a'_t:\,d_t(a'_t,h_t)=a_t}
  g_t(a'_t\mid h_t).
\]

\begin{assumption}[Baseline conditioning event]
\label{assump:baseline-v}
For discrete \(V\), assume
\[
  p_v = \mathbb{P}(V=v)>0
\]
for every subgroup value \(v\) under consideration. For continuous \(V\),
\(\theta_{\mathbbm d}(v)\) is interpreted as a version of
\(\mathbb{E}_{P^{\mathbbm d}}(Y\mid V=v)\), defined for \(v\) in the support of
\(V\).
\end{assumption}

Assumptions~\ref{assump:consistency} and~\ref{assump:seq-randomization} give the
causal interpretation of the post-policy law. Assumption~\ref{assump:lmtp-support}
ensures that the policy does not assign treatment values outside the observed
conditional support. Assumption~\ref{assump:policy-regularity} ensures that the
post-policy treatment law \(g_t^{\mathbbm d}\) and density ratios used for
estimation are well-defined. Assumption~\ref{assump:baseline-v} ensures that the
conditional target is defined for the subgroup or covariate values of interest.

\subsection{Iterated conditional mean representation}
\label{subsec:iterated-mean-representation}

Define the terminal outcome regression
\[
  m_{\tau}^{\mathbbm d}(a_\tau,h_\tau)
  =
  \mathbb{E}_{P}
  \left[
    Y
    \mid
    A_\tau=a_\tau,H_\tau=h_\tau
  \right].
\]
For \(t=\tau-1,\ldots,0\), define the sequential regressions recursively by
\[
  m_t^{\mathbbm d}(a_t,h_t)
  =
  \mathbb{E}_{P}
  \left[
    m_{t+1}^{\mathbbm d}
    \{A_{t+1}^{\mathbbm d},H_{t+1}\}
    \,\middle|\,
    A_t=a_t,H_t=h_t
  \right].
\]
Here
\[
  A_{t+1}^{\mathbbm d}
  =
  d_{t+1}(A_{t+1},H_{t+1})
\]
is the treatment value assigned by the modified treatment policy at time
\(t+1\), evaluated using the observed history \(H_{t+1}\). This recursion is the
iterated conditional mean version of the longitudinal \(g\)-formula.

With this notation, the marginal LMTP mean is identified as
\[
  \psi_{\mathbbm d}
  =
  \mathbb{E}_{P}
  \left[
    m_0^{\mathbbm d}(A_0^{\mathbbm d},L_0)
  \right],
\]
and the conditional LMTP mean is identified as
\[
  \theta_{\mathbbm d}(v)
  =
  \mathbb{E}_{P}
  \left[
    m_0^{\mathbbm d}(A_0^{\mathbbm d},L_0)
    \mid V=v
  \right].
\]

\begin{theorem}[Iterated conditional mean identification]
\label{thm:iterated-mean-identification}
Under the assumptions~\ref{assump:consistency}--\ref{assump:baseline-v},
\begin{equation} \label{eqn:icm1}
  \theta_{\mathbbm d}(v)
  =
  \mathbb{E}_{P}
  \left[
    m_0^{\mathbbm d}(A_0^{\mathbbm d},L_0)
    \mid V=v
  \right].
\end{equation}
For two policies \(\mathbbm d_1\) and \(\mathbbm d_0\),
\[
  \tau(v)
  =
  \mathbb{E}_{P}
  \left[
    m_0^{\mathbbm d_1}(A_0^{\mathbbm d_1},L_0)
    -
    m_0^{\mathbbm d_0}(A_0^{\mathbbm d_0},L_0)
    \mid V=v
  \right].
\]
\end{theorem}

\begin{proof}
The result follows by repeated application of the law of iterated expectations. At the final
time point,
\[
 m_\tau^{\mathbbm d}(a_\tau,h_\tau)
  =
  \mathbb{E}_{P}(Y\mid A_\tau=a_\tau,H_\tau=h_\tau)
\]
is the conditional mean of the outcome given the final treatment and history.
At time \(\tau-1\), the regression
\[
  m_{\tau-1}^{\mathbbm d}(a_{\tau-1},h_{\tau-1})
  =
  \mathbb{E}_{P}
  \left[
    m_{\tau}^{\mathbbm d}(A_\tau^{\mathbbm d},H_\tau)
    \mid A_{\tau-1}=a_{\tau-1},H_{\tau-1}=h_{\tau-1}
  \right]
\]
integrates the final outcome regression over the conditional law of the next
covariate and treatment, with the next treatment evaluated under the modified
policy. Repeating this argument backward through
\(t=\tau-2,\ldots,0\) gives the \(g\)-formula mean under
\(P^{\mathbbm d}\). Conditioning on \(V=v\) gives \ref{eqn:icm1}. The contrast
result follows by applying the same argument to each policy and taking the
difference.
\end{proof}

\subsection{Uncentered influence-function pseudo-outcomes}
\label{subsec:uncentered-eif-pseudooutcomes}

The preceding representations identify the conditional LMTP mean. For
estimation, it is useful to construct a pseudo-outcome whose conditional
expectation is the target function \(v\mapsto \theta_{\mathbbm d}(v)\). This
pseudo-outcome is obtained from the efficient influence function for the marginal
LMTP mean by adding back the marginal mean.

Let
\[
  \eta^{\mathbbm d}
  =
  \left(
    m_0^{\mathbbm d},\ldots,m_\tau^{\mathbbm d},
    r_0^{\mathbbm d},\ldots,r_\tau^{\mathbbm d}
  \right)
\]
denote the nuisance functions for policy \(\mathbbm d\). Define
\[
  \omega_s^{\mathbbm d}(O)
  =
  \prod_{k=0}^{s}
  r_k^{\mathbbm d}(A_k,H_k),
  \qquad
  s=0,\ldots,\tau.
\]
For \(s=0,\ldots,\tau-1\), define the sequential residual
\[
  \Delta_s^{\mathbbm d}(O;\eta^{\mathbbm d})
  =
  m_{s+1}^{\mathbbm d}
  \{A_{s+1}^{\mathbbm d},H_{s+1}\}
  -
  m_s^{\mathbbm d}(A_s,H_s),
\]
and define the terminal residual
\[
  \Delta_\tau^{\mathbbm d}(O;\eta^{\mathbbm d})
  =
  Y - m_\tau^{\mathbbm d}(A_\tau,H_\tau).
\]
Thus each residual compares the next-stage pseudo-outcome to its conditional
mean given the current treatment and history.

The uncentered influence-function pseudo-outcome is
\[
  \varphi_{\mathbbm d}(O;\eta^{\mathbbm d})
  =
  m_0^{\mathbbm d}(A_0^{\mathbbm d},L_0)
  +
  \sum_{s=0}^{\tau}
  \omega_s^{\mathbbm d}(O)
  \Delta_s^{\mathbbm d}(O;\eta^{\mathbbm d}).
\]
Equivalently,
\begin{align*}
  \varphi_{\mathbbm d}(O;\eta^{\mathbbm d})
  =
  m_0^{\mathbbm d}(A_0^{\mathbbm d},L_0)
 +
  \sum_{s=0}^{\tau-1}
  \left\{
    \prod_{k=0}^{s}
    r_k^{\mathbbm d}(A_k,H_k)
  \right\}
  \Big[
    m_{s+1}^{\mathbbm d}
    \{ & A_{s+1}^{\mathbbm d},H_{s+1}\}
    -
    m_s^{\mathbbm d}(A_s,H_s)
  \Big] \\ 
  & +
  \left\{
    \prod_{k=0}^{\tau}
    r_k^{\mathbbm d}(A_k,H_k)
  \right\}
  \left[
    Y - m_\tau^{\mathbbm d}(A_\tau,H_\tau)
  \right].
\end{align*}
The product in this expression is the product of density ratios up to time \(s\);
the terms in square brackets are sequential regression residuals. Thus the EIF
is a sum of weighted residuals.

At the true nuisance functions, the centered efficient influence function for
the marginal LMTP mean \(\psi_{\mathbbm d}\) is
\[
  D_{\mathbbm d}(O)
  =
  \varphi_{\mathbbm d}(O;\eta^{\mathbbm d})
  -
  \psi_{\mathbbm d}.
\]
Equivalently,
\[
  \varphi_{\mathbbm d}(O;\eta^{\mathbbm d})
  =
  \psi_{\mathbbm d}
  +
  D_{\mathbbm d}(O).
\]
Therefore
\[
  \mathbb{E}_{P}
  \{
    \varphi_{\mathbbm d}(O;\eta^{\mathbbm d})
  \}
  =
  \psi_{\mathbbm d}.
\]

The key property for conditional effect estimation is that, since
\(V\subseteq L_0\), the conditional mean of the uncentered pseudo-outcome is the
conditional LMTP mean.

\begin{theorem}[Pseudo-outcome representation of the conditional LMTP mean]
\label{thm:pseudooutcome-identification}
Under Assumptions~\ref{assump:consistency}--\ref{assump:baseline-v},
\[
  \mathbb{E}_{P}
  \left[
    \varphi_{\mathbbm d}(O;\eta^{\mathbbm d})
    \mid V=v
  \right]
  =
  \theta_{\mathbbm d}(v).
\]
Therefore, for two policies \(\mathbbm d_1\) and \(\mathbbm d_0\), if
\[
  \varphi_{\tau}(O;\eta)
  =
  \varphi_{\mathbbm d_1}(O;\eta^{\mathbbm d_1})
  -
  \varphi_{\mathbbm d_0}(O;\eta^{\mathbbm d_0}),
\]
then
\[
  \mathbb{E}_{P}
  \left[
    \varphi_{\tau}(O;\eta)
    \mid V=v
  \right]
  =
  \tau(v).
\]
\end{theorem}

\begin{proof}
By Theorem~\ref{thm:iterated-mean-identification},
\[
  \mathbb{E}_{P}
  \left[
    m_0^{\mathbbm d}(A_0^{\mathbbm d},L_0)
    \mid V=v
  \right]
  =
  \theta_{\mathbbm d}(v).
\]
It remains to show that the weighted residual terms have conditional mean zero
given \(V\).

Fix \(s\in\{0,\ldots,\tau-1\}\). Since
\(\omega_s^{\mathbbm d}(O)\) is measurable with respect to
\((A_s,H_s)\), and since \(V\subseteq L_0\subseteq H_s\), iterated expectation
gives
\[
  \mathbb{E}_{P}
  \left[
    \omega_s^{\mathbbm d}(O)
    \Delta_s^{\mathbbm d}(O;\eta^{\mathbbm d})
    \mid V
  \right]
  =
  \mathbb{E}_{P}
  \left[
    \omega_s^{\mathbbm d}(O)
    \mathbb{E}_{P}
    \left\{
      \Delta_s^{\mathbbm d}(O;\eta^{\mathbbm d})
      \mid A_s,H_s
    \right\}
    \mid V
  \right].
\]
By the definition of \(m_s^{\mathbbm d}\),
\[
  \mathbb{E}_{P}
  \left[
    m_{s+1}^{\mathbbm d}
    \{A_{s+1}^{\mathbbm d},H_{s+1}\}
    \mid A_s,H_s
  \right]
  =
  m_s^{\mathbbm d}(A_s,H_s),
\]
so
\[
  \mathbb{E}_{P}
  \left\{
    \Delta_s^{\mathbbm d}(O;\eta^{\mathbbm d})
    \mid A_s,H_s
  \right\}
  =
  0.
\]
Therefore
\[
  \mathbb{E}_{P}
  \left[
    \omega_s^{\mathbbm d}(O)
    \Delta_s^{\mathbbm d}(O;\eta^{\mathbbm d})
    \mid V
  \right]
  =
  0.
\]

For the terminal term,
\[
  \mathbb{E}_{P}
  \left[
    Y - m_\tau^{\mathbbm d}(A_\tau,H_\tau)
    \mid A_\tau,H_\tau
  \right]
  =
  0
\]
by the definition of \(m_\tau^{\mathbbm d}\). The same iterated-expectation
argument gives
\[
  \mathbb{E}_{P}
  \left[
    \omega_\tau^{\mathbbm d}(O)
    \Delta_\tau^{\mathbbm d}(O;\eta^{\mathbbm d})
    \mid V
  \right]
  =
  0.
\]
Thus all weighted residual terms in
\(\varphi_{\mathbbm d}(O;\eta^{\mathbbm d})\) have conditional mean zero given
\(V\). Hence
\[
  \mathbb{E}_{P}
  \left[
    \varphi_{\mathbbm d}(O;\eta^{\mathbbm d})
    \mid V=v
  \right]
  =
  \mathbb{E}_{P}
  \left[
    m_0^{\mathbbm d}(A_0^{\mathbbm d},L_0)
    \mid V=v
  \right]
  =
  \theta_{\mathbbm d}(v).
\]
The contrast result follows by linearity of conditional expectation.
\end{proof}

\begin{remark}[Role of the pseudo-outcome]
\label{rem:pseudooutcome-role}
The pseudo-outcome
\(\varphi_{\mathbbm d}(O;\eta^{\mathbbm d})\) is not a new causal parameter. It
is a representation of the same identified LMTP mean. Its role is statistical:
it converts estimation of the conditional LMTP mean into a regression problem.
For a general conditioning variable \(V\), the DR-Learner estimates
\[
  v
  \mapsto
  \mathbb{E}
  \left[
    \varphi_{\mathbbm d}(O;\eta^{\mathbbm d})
    \mid V=v
  \right]
\]
by replacing the unknown nuisance functions with estimates and then regressing
the resulting estimated pseudo-outcomes on \(V\).
\end{remark}

\section{Estimation}
\label{sec:estimation}

We now describe estimators for the conditional LMTP parameters introduced in
Section~\ref{sec:setup}. We consider two complementary estimation strategies.
First, when \(V\) is discrete with finite support \(\mathcal V\), we construct a
targeted minimum loss-based estimator of the finite-dimensional vector
\[
  \{\theta_{\mathbbm d}(v):v\in\mathcal V\}.
\]
This estimator targets the efficient influence-function estimating equations
for all subgroups jointly. Second, when \(V\) is continuous, high-dimensional, or
when the goal is to estimate a smooth conditional effect function, we construct
a DR-Learner by regressing estimated uncentered influence-function
pseudo-outcomes on \(V\).

Throughout this section, let
\[
  \eta^{\mathbbm d}
  =
  \left(
    m_0^{\mathbbm d},\ldots,m_\tau^{\mathbbm d},
    r_0^{\mathbbm d},\ldots,r_\tau^{\mathbbm d}
  \right)
\]
denote the nuisance functions for policy \(\mathbbm d\). For two policies
\(\mathbbm d_1\) and \(\mathbbm d_0\), all steps below may be applied separately
to each policy and then contrasted.

\subsection{Cross-fitting and nuisance estimation}
\label{subsec:crossfitting-nuisance}

Flexible estimation of the nuisance functions is essential in longitudinal
settings, where the treatment mechanism and sequential outcome regressions may
depend on high-dimensional histories. We use cross-fitting to separate nuisance
estimation from evaluation of the estimating equations. This allows the nuisance
functions to be estimated with flexible regression or machine-learning methods
without requiring empirical process restrictions such as Donsker conditions.

Let
\[
  \mathcal I_1,\ldots,\mathcal I_J
\]
be a partition of \(\{1,\ldots,n\}\) into \(J\) folds. For each fold \(j\), let
\(\mathcal I_{-j}\) denote its complement. Nuisance functions are fit using only
the training sample \(\mathcal I_{-j}\), and predictions are evaluated only on
the validation sample \(\mathcal I_j\). For observation \(i\), let \(j(i)\)
denote the fold containing \(i\).

For each policy \(\mathbbm d\), each fold \(j\), and each time point
\(t=0,\ldots,\tau\), estimate the density ratio
\[
  r_t^{\mathbbm d}(a_t,h_t)
  =
  \frac{
    g_t^{\mathbbm d}(a_t\mid h_t)
  }{
    g_t(a_t\mid h_t)
  }
\]
using the training sample \(\mathcal I_{-j}\). Denote the resulting
cross-fitted prediction for observation \(i\) by
\[
  \widehat r_{t,i}^{\mathbbm d}
  =
  \widehat r_{t,j(i)}^{\mathbbm d}(A_{t,i},H_{t,i}).
\]
The estimated cumulative density ratio through time \(s\) is
\[
  \widehat\omega_{s,i}^{\mathbbm d}
  =
  \prod_{t=0}^{s}
  \widehat r_{t,i}^{\mathbbm d}.
\]

The sequential outcome regressions are also estimated recursively using
cross-fitting. Starting at the final time point, estimate
\[
  m_\tau^{\mathbbm d}(a_\tau,h_\tau)
  =
  \mathbb{E}(Y\mid A_\tau=a_\tau,H_\tau=h_\tau)
\]
on \(\mathcal I_{-j}\), and evaluate it on \(\mathcal I_j\) at both the observed
and policy-modified treatment values as required. For
\(t=\tau-1,\ldots,0\), estimate
\[
  m_t^{\mathbbm d}(a_t,h_t)
  =
  \mathbb{E}
  \left[
    m_{t+1}^{\mathbbm d}
    \{A_{t+1}^{\mathbbm d},H_{t+1}\}
    \mid A_t=a_t,H_t=h_t
  \right].
\]
In practice, the regression outcome at time \(t\) is the previously estimated
next-stage pseudo-outcome
\[
  \widehat m_{t+1,j}^{\mathbbm d}
  \{A_{t+1}^{\mathbbm d},H_{t+1}\},
\]
constructed using the training sample. The resulting cross-fitted predictions
for observation \(i\) are denoted by
\[
  \widehat m_{t,i}^{\mathbbm d}(A_{t,i},H_{t,i})
  \qquad\text{and}\qquad
  \widehat m_{t,i}^{\mathbbm d}(A_{t,i}^{\mathbbm d},H_{t,i}).
\]

Given these nuisance estimates, define the cross-fitted estimated residuals
\[
  \widehat\Delta_{s,i}^{\mathbbm d}
  =
  \widehat m_{s+1,i}^{\mathbbm d}
  \{A_{s+1,i}^{\mathbbm d},H_{s+1,i}\}
  -
  \widehat m_{s,i}^{\mathbbm d}(A_{s,i},H_{s,i}),
  \qquad s=0,\ldots,\tau-1,
\]
and
\[
  \widehat\Delta_{\tau,i}^{\mathbbm d}
  =
  Y_i
  -
  \widehat m_{\tau,i}^{\mathbbm d}(A_{\tau,i},H_{\tau,i}).
\]
The estimated uncentered influence-function pseudo-outcome is then
\[
  \widehat\varphi_{\mathbbm d,i}
  =
  \widehat m_{0,i}^{\mathbbm d}(A_{0,i}^{\mathbbm d},L_{0,i})
  +
  \sum_{s=0}^{\tau}
  \widehat\omega_{s,i}^{\mathbbm d}
  \widehat\Delta_{s,i}^{\mathbbm d}.
\]
For a contrast of two policies, define
\[
  \widehat\varphi_{\tau,i}
  =
  \widehat\varphi_{\mathbbm d_1,i}
  -
  \widehat\varphi_{\mathbbm d_0,i}.
\]

\begin{remark}
The notation \(\widehat\varphi_{\tau,i}\) denotes the estimated pseudo-outcome
for the conditional contrast \(\tau(v)\), not a time-indexed quantity. When
there is risk of confusion, we write
\(\widehat\varphi_{\mathbbm d_1-\mathbbm d_0,i}\).
\end{remark}

\subsection{Targeted estimation for finite baseline subgroups}
\label{subsec:subgroup-tmle}

Suppose \(V\) is discrete with finite support
\[
  \mathcal V = \{v_1,\ldots,v_K\},
\]
and let
\[
  p_v = \mathbb{P}(V=v).
\]
The conditional LMTP mean for subgroup \(v\) can be written as
\[
  \theta_{\mathbbm d}(v)
  =
  \frac{
    \mathbb{E}
    \left[
      \mathds{1}(V=v)
      m_0^{\mathbbm d}(A_0^{\mathbbm d},L_0)
    \right]
  }{
    p_v
  }.
\]
The efficient influence function for \(\theta_{\mathbbm d}(v)\) is
\[
  D_{\mathbbm d,v}(O)
  =
  \frac{\mathds{1}(V=v)}{p_v}
  \left\{
    \varphi_{\mathbbm d}(O;\eta^{\mathbbm d})
    -
    \theta_{\mathbbm d}(v)
  \right\}.
\]
Equivalently,
\[
  D_{\mathbbm d,v}(O)
  =
  \frac{\mathds{1}(V=v)}{p_v}
  \left[
    m_0^{\mathbbm d}(A_0^{\mathbbm d},L_0)
    -
    \theta_{\mathbbm d}(v)
    +
    \sum_{s=0}^{\tau}
    \omega_s^{\mathbbm d}(O)
    \Delta_s^{\mathbbm d}(O;\eta^{\mathbbm d})
  \right].
\]
For a contrast
\[
  \tau(v)
  =
  \theta_{\mathbbm d_1}(v)-\theta_{\mathbbm d_0}(v),
\]
the efficient influence function is
\[
  D_{\tau,v}(O)
  =
  D_{\mathbbm d_1,v}(O)
  -
  D_{\mathbbm d_0,v}(O).
\]

The subgroup TMLE targets the sequential outcome regressions so that the
empirical mean of the estimated efficient influence function is approximately
zero for each \(v\in\mathcal V\). Let
\[
  \widehat p_v
  =
  \mathbb{P}_n\{\mathds{1}(V=v)\}.
\]
For \(s=0,\ldots,\tau\), define the subgroup-specific clever covariate
\[
  \widehat H_{s,v}^{\mathbbm d}(O_i)
  =
  \frac{\mathds{1}(V_i=v)}{\widehat p_v}
  \widehat\omega_{s,i}^{\mathbbm d}.
\]
Collecting these across \(v\in\mathcal V\) gives the vector clever covariate
\[
  \widehat H_s^{\mathbbm d}(O_i)
  =
  \left(
    \widehat H_{s,v_1}^{\mathbbm d}(O_i),
    \ldots,
    \widehat H_{s,v_K}^{\mathbbm d}(O_i)
  \right)^\top.
\]

The targeting step proceeds backward in time. Let
\[
  \widehat Y_{\tau+1,i}^{\mathbbm d,*}=Y_i.
\]
For \(s=\tau,\tau-1,\ldots,0\), fit a fluctuation submodel for the current
outcome regression using \(\widehat Y_{s+1,i}^{\mathbbm d,*}\) as the outcome,
\(\widehat m_{s,i}^{\mathbbm d}(A_{s,i},H_{s,i})\) as the initial prediction, and
\(\widehat H_s^{\mathbbm d}(O_i)\) as the clever covariate. With an identity
link, the fluctuation model is
\[
  \widehat m_{s,\epsilon}^{\mathbbm d}(A_{s,i},H_{s,i})
  =
  \widehat m_{s}^{\mathbbm d}(A_{s,i},H_{s,i})
  +
  \epsilon_s^\top
  \widehat H_s^{\mathbbm d}(O_i),
\]
where \(\epsilon_s\in\mathbb{R}^K\). With a bounded outcome and logit link, the
fluctuation model is
\[
  \operatorname{logit}
  \widehat m_{s,\epsilon}^{\mathbbm d}(A_{s,i},H_{s,i})
  =
  \operatorname{logit}
  \widehat m_{s}^{\mathbbm d}(A_{s,i},H_{s,i})
  +
  \epsilon_s^\top
  \widehat H_s^{\mathbbm d}(O_i).
\]
Let \(\widehat\epsilon_s\) denote the fitted fluctuation coefficient and write
\[
  \widehat m_{s}^{\mathbbm d,*}
  =
  \widehat m_{s,\widehat\epsilon_s}^{\mathbbm d}.
\]

After updating the regression at time \(s\), define the next outcome for the
backward recursion by evaluating the updated regression under the modified
treatment value:
\[
  \widehat Y_{s,i}^{\mathbbm d,*}
  =
  \widehat m_{s}^{\mathbbm d,*}
  (A_{s,i}^{\mathbbm d},H_{s,i}).
\]
This updated pseudo-outcome is then used in the fluctuation at time \(s-1\).

The subgroup TMLE of the policy-specific conditional mean is the plug-in
estimator
\[
  \widehat\theta_{\mathbbm d,\mathrm{tmle}}(v)
  =
  \mathbb{P}_n
  \left[
    \frac{\mathds{1}(V=v)}{\widehat p_v}
    \widehat m_0^{\mathbbm d,*}(A_0^{\mathbbm d},L_0)
  \right].
\]
For two policies \(\mathbbm d_1\) and \(\mathbbm d_0\), the contrast estimator is
\[
  \widehat\tau_{\mathrm{tmle}}(v)
  =
  \widehat\theta_{\mathbbm d_1,\mathrm{tmle}}(v)
  -
  \widehat\theta_{\mathbbm d_0,\mathrm{tmle}}(v).
\]

The estimated efficient influence function for the policy-specific subgroup
mean is
\[
  \widehat D_{\mathbbm d,v}(O_i)
  =
  \frac{\mathds{1}(V_i=v)}{\widehat p_v}
  \left[
    \widehat m_0^{\mathbbm d,*}(A_{0,i}^{\mathbbm d},L_{0,i})
    -
    \widehat\theta_{\mathbbm d,\mathrm{tmle}}(v)
    +
    \sum_{s=0}^{\tau}
    \widehat\omega_{s,i}^{\mathbbm d}
    \widehat\Delta_{s,i}^{\mathbbm d,*}
  \right],
\]
where
\[
  \widehat\Delta_{s,i}^{\mathbbm d,*}
  =
  \widehat m_{s+1}^{\mathbbm d,*}
  (A_{s+1,i}^{\mathbbm d},H_{s+1,i})
  -
  \widehat m_{s}^{\mathbbm d,*}(A_{s,i},H_{s,i}),
  \qquad s=0,\ldots,\tau-1,
\]
and
\[
  \widehat\Delta_{\tau,i}^{\mathbbm d,*}
  =
  Y_i
  -
  \widehat m_{\tau}^{\mathbbm d,*}(A_{\tau,i},H_{\tau,i}).
\]
For a contrast, define
\[
  \widehat D_{\tau,v}(O_i)
  =
  \widehat D_{\mathbbm d_1,v}(O_i)
  -
  \widehat D_{\mathbbm d_0,v}(O_i).
\]

For a single subgroup \(v\), the estimated standard error of
\(\widehat\tau_{\mathrm{tmle}}(v)\) is
\[
  \widehat{\mathrm{se}}\{\widehat\tau_{\mathrm{tmle}}(v)\}
  =
  \left[
    \frac{1}{n}
    \mathbb{P}_n
    \left\{
      \widehat D_{\tau,v}(O)^2
    \right\}
  \right]^{1/2}.
\]
For simultaneous inference over \(v\in\mathcal V\), define the empirical
covariance matrix
\[
  \widehat\Sigma_{\tau}
  =
  \mathbb{P}_n
  \left[
    \widehat D_{\tau}(O)
    \widehat D_{\tau}(O)^\top
  \right],
\]
where
\[
  \widehat D_{\tau}(O)
  =
  \left(
    \widehat D_{\tau,v_1}(O),
    \ldots,
    \widehat D_{\tau,v_K}(O)
  \right)^\top.
\]
Then
\[
  \sqrt n
  \left\{
    \widehat\tau_{\mathrm{tmle}}-\tau
  \right\}
\]
is approximated by a mean-zero multivariate normal distribution with covariance
matrix \(\widehat\Sigma_\tau\). Simultaneous confidence intervals may be formed
using the quantiles of
\[
  \max_{v\in\mathcal V}
  \left|
    Z_v / \widehat\sigma_v
  \right|,
  \qquad
  Z\sim N(0,\widehat\Sigma_\tau),
\]
where \(\widehat\sigma_v^2\) is the \(v\)th diagonal element of
\(\widehat\Sigma_\tau\).

\begin{remark}
The targeting step above updates the full vector of subgroup parameters jointly
at each time point. This is preferable to fitting a separate TMLE within each
subgroup when some subgroups are small, because the nuisance regressions are
estimated using the full sample while the clever covariates target the
subgroup-specific components of the efficient influence function.
\end{remark}

\subsection{DR-Learner for conditional LMTP effects}
\label{subsec:dr-learner}

When \(V\) is continuous or high-dimensional, the target
\[
  v\mapsto \theta_{\mathbbm d}(v)
\]
is a conditional mean function rather than a finite-dimensional parameter. The
pseudo-outcome representation in
Theorem~\ref{thm:pseudooutcome-identification} suggests a direct meta-learning
approach: estimate the uncentered influence-function pseudo-outcome and regress
it on \(V\).

For a single policy \(\mathbbm d\), the DR-Learner proceeds as follows.

\begin{enumerate}
  \item Split the sample into folds
  \(\mathcal I_1,\ldots,\mathcal I_J\).

  \item On each training sample \(\mathcal I_{-j}\), estimate the nuisance
  functions
  \[
    \eta^{\mathbbm d}
    =
    (m_0^{\mathbbm d},\ldots,m_\tau^{\mathbbm d},
     r_0^{\mathbbm d},\ldots,r_\tau^{\mathbbm d}).
  \]

  \item For each observation \(i\in\mathcal I_j\), construct the cross-fitted
  pseudo-outcome
  \[
    \widehat\varphi_{\mathbbm d,i}
    =
    \widehat m_{0,i}^{\mathbbm d}(A_{0,i}^{\mathbbm d},L_{0,i})
    +
    \sum_{s=0}^{\tau}
    \widehat\omega_{s,i}^{\mathbbm d}
    \widehat\Delta_{s,i}^{\mathbbm d}.
  \]

  \item Regress \(\widehat\varphi_{\mathbbm d,i}\) on \(V_i\) using a chosen
  second-stage regression method. The fitted regression function is
  \[
    \widehat\theta_{\mathbbm d,\mathrm{dr}}(v)
    =
    \widehat{\mathbb{E}}_n
    \left[
      \widehat\varphi_{\mathbbm d}
      \mid V=v
    \right].
  \]
\end{enumerate}

For two policies \(\mathbbm d_1\) and \(\mathbbm d_0\), construct the contrast
pseudo-outcome
\[
  \widehat\varphi_{\tau,i}
  =
  \widehat\varphi_{\mathbbm d_1,i}
  -
  \widehat\varphi_{\mathbbm d_0,i},
\]
and regress it on \(V_i\):
\[
  \widehat\tau_{\mathrm{dr}}(v)
  =
  \widehat{\mathbb{E}}_n
  \left[
    \widehat\varphi_{\tau}
    \mid V=v
  \right].
\]
Equivalently,
\[
  \widehat\tau_{\mathrm{dr}}(v)
  =
  \widehat\theta_{\mathbbm d_1,\mathrm{dr}}(v)
  -
  \widehat\theta_{\mathbbm d_0,\mathrm{dr}}(v),
\]
if the same second-stage regression method is applied separately to the two
policy-specific pseudo-outcomes.

The second-stage regression may be chosen to reflect the scientific structure
expected in the conditional effect. If \(V\) is low-dimensional and continuous,
local linear regression, splines, generalized additive models, or local
polynomial regression are natural choices. If \(V\) is multivariate, random
forests, gradient boosting, neural networks, or Super Learner ensembles may be
used. If the scientific question concerns a pre-specified low-dimensional
projection of the conditional effect, one may regress the pseudo-outcome on a
parametric or semiparametric working model for \(V\).

The corresponding benchmark estimator is the regression that would be obtained
if the nuisance functions were known:
\[
  \widetilde\theta_{\mathbbm d}(v)
  =
  \widehat{\mathbb{E}}_n
  \left[
    \varphi_{\mathbbm d}(O;\eta^{\mathbbm d})
    \mid V=v
  \right],
\]
and, for a contrast,
\[
  \widetilde\tau(v)
  =
  \widehat{\mathbb{E}}_n
  \left[
    \varphi_{\tau}(O;\eta)
    \mid V=v
  \right].
\]
The theoretical results in Section~\ref{sec:theory} compare the feasible
DR-Learner to this benchmark regression. In particular, oracle equivalence holds
when the smoothed second-order bias induced by nuisance estimation is negligible
relative to the oracle second-stage regression error.

\begin{remark}
For finite discrete \(V\), the DR-Learner with a saturated second-stage
regression is simply the subgroup mean of the estimated pseudo-outcomes. Thus
the finite-subgroup estimator can be viewed as the discrete special case of the
pseudo-outcome regression strategy. The targeted estimator, however, has the
additional advantage that it updates the nuisance regressions to solve the
efficient influence-function equations for the subgroup parameters.
\end{remark}

\begin{remark}[Second-stage regression and HAL]
Uncertainty quantification for the DR-Learner depends on both the
second-stage regression method and the target of inference. For
finite-dimensional projections of the conditional effect, standard asymptotic
linearity can often be obtained from the corresponding projected estimating
equations. For pointwise or uniform inference on a nonparametric conditional
effect curve, additional analysis is required.

One strategically attractive choice for the second-stage regression is the
highly adaptive lasso (HAL). HAL estimates cadlag functions with finite
sectional variation norm by empirical risk minimization over a large dictionary
of tensor-product indicator basis functions subject to an \(\ell_1\) constraint.
For regression-type function estimation, HAL attains loss-based convergence
rates that typically imply \(L_2(P_0)\) convergence at rate
\(n^{-1/3}\), up to logarithmic factors, under only a finite sectional variation
norm condition. This rate is faster than \(n^{-1/4}\), which is often the
critical threshold appearing in second-order remainder conditions for
doubly robust and targeted estimators. HAL is therefore appealing as a
second-stage learner because it can estimate flexible conditional effect
functions while retaining nonparametric rate guarantees under relatively weak
global regularity assumptions.

Undersmoothed HAL may be particularly relevant if the goal is inference rather
than only prediction. Existing work shows that undersmoothed HAL can yield
asymptotically efficient plug-in estimators of pathwise differentiable features
of an infinite-dimensional regression function, and higher-order spline
extensions of HAL provide pointwise asymptotic normality and uniform
convergence rates for function estimation under stronger smoothness conditions.
These results suggest that HAL-based second-stage regression may provide a
promising route toward inference for conditional LMTP effect curves. However,
establishing pointwise or uniform confidence bands for the LMTP DR-Learner
requires theory specific to the pseudo-outcome, the second-stage target, and the
chosen smoothing or undersmoothing rule. We therefore establish a general
oracle-equivalence result for the DR-Learner in
Section~\ref{sec:theory} and leave method-specific pointwise and uniform
inference for conditional LMTP effect curves to future work.
\end{remark}


\subsection{Projections onto working marginal structural models}
\label{subsec:msm-projection}

The subgroup TMLE of Section~\ref{subsec:subgroup-tmle} and the DR-Learner of
Section~\ref{subsec:dr-learner} occupy two extremes: fully saturated inference
over a finite \(\mathcal V\), and fully nonparametric estimation of
\(v\mapsto\theta_{\d}(v)\) at nonparametric rates. Applied reporting often
calls for something in between --- a low-dimensional, interpretable summary of
how the policy effect varies with \(V\), with root-\(n\) inference --- or for
pointwise and uniform inference on the conditional effect curve itself. Both
are obtained by regressing the influence-function pseudo-outcome on a basis in
\(V\): a \emph{fixed} basis yields a working marginal structural model (MSM)
whose projection coefficients are pathwise differentiable, while a basis of
\emph{growing} dimension, undersmoothed so that approximation bias is
negligible, yields valid inference on \(\theta_{\d}\) itself. We present the
two regimes together because they share one estimator.

\subsection{Fixed-dimensional working models}
\label{subsubsec:msm-fixed}

Let \(b(v)=\{b_1(v),\ldots,b_k(v)\}^\top\) be a fixed, user-specified vector of
basis functions --- for example \(b(v)=(1,v)^\top\) for a linear trend in a
scalar modifier, or a coarse spline basis --- and let \(h(v)\ge 0\) be a
user-specified weight function, with \(h\equiv 1\) as the default. Define the
working-model projection of the conditional LMTP mean,
\[
  \beta_{\d}
  =
  \arg\min_{\beta\in\R^k}
  \E
  \left[
    h(V)
    \left\{
      \theta_{\d}(V)
      -
      b(V)^\top\beta
    \right\}^2
  \right]
  =
  M^{-1}
  \,
  \E
  \left\{
    h(V)\,b(V)\,\theta_{\d}(V)
  \right\},
  \qquad
  M
  =
  \E
  \left\{
    h(V)\,b(V)\,b(V)^\top
  \right\},
\]
assuming \(M\) is nonsingular. The parameter \(\beta_{\d}\) is defined without
assuming the working model is correct: \(v\mapsto b(v)^\top\beta_{\d}\) is the
weighted least-squares projection of \(\theta_{\d}\) onto the span of \(b\),
and retains this interpretation under arbitrary misspecification
\cite{neugebauer2007nonparametric,rosenblum2010msm}. When the working model
happens to hold, \(b(v)^\top\beta_{\d}=\theta_{\d}(v)\). For a contrast of two
policies, \(\beta_\tau=\beta_{\d_1}-\beta_{\d_0}\) is the projection of
\(\tau(\cdot)\), obtained by replacing \(\theta_{\d}\) with \(\tau\) throughout,
or equivalently by using the contrast pseudo-outcome
\(\varphi_\tau=\varphi_{\d_1}-\varphi_{\d_0}\).

Because
\(\E\{\varphi_{\d}(O;\eta_0^{\d})\mid V=v\}=\theta_{\d}(v)\) by
Theorem~\ref{thm:pseudooutcome-identification}, the projection coefficients
admit the moment representation
\(\beta_{\d}=M^{-1}\E\{h(V)b(V)\varphi_{\d}(O;\eta_0^{\d})\}\), which is a
smooth finite-dimensional functional of the observed data law.

\begin{proposition}[Efficient influence function of the projection]
\label{prop:msm-eif}
Under Assumptions~\ref{assump:consistency}--\ref{assump:baseline-v} and
nonsingular \(M\), the parameter \(\beta_{\d}\) is pathwise differentiable with
efficient influence function
\[
  D_{\beta,\d}(O)
  =
  M^{-1}
  \,
  h(V)\,b(V)
  \left\{
    \varphi_{\d}(O;\eta_0^{\d})
    -
    b(V)^\top\beta_{\d}
  \right\}.
\]
For the contrast, \(D_{\beta,\tau}=D_{\beta,\d_1}-D_{\beta,\d_0}\).
\end{proposition}

\begin{proof}
The parameter solves the moment equation
\(\E\,U(O;\beta)=0\) with
\(U(O;\beta)=h(V)b(V)\{\varphi_{\d}(O;\eta_0^{\d})-b(V)^\top\beta\}\). The
functional \(P\mapsto\E_P\{h(V)b(V)\theta_{\d,P}(V)\}\) has efficient influence
function \(h(V)b(V)\varphi_{\d}(O;\eta_0^{\d})\) minus its mean, since
\(\varphi_{\d}\) is the uncentered efficient influence function for the LMTP
mean and \(h\,b\) is a fixed baseline-measurable weight; the contribution of
the marginal law of \(V\) through \(M\) is absorbed by centering at
\(b(V)^\top\beta_{\d}\), which renders \(D_{\beta,\d}\) mean zero by the normal
equations. The delta method for estimating equations, with derivative matrix
\(\partial_\beta\,\E\,U=-M\), gives the stated form. Details are provided in
the Supplementary Material.
\end{proof}

The estimator is the weighted least-squares regression of the cross-fitted
pseudo-outcomes of Section~\ref{subsec:crossfitting-nuisance} on the basis:
\[
  \widehat\beta_{\d}
  =
  \widehat M^{-1}
  \,
  \mathbb{P}_n
  \left\{
    h(V)\,b(V)\,
    \widehat\varphi_{\d}(O)
  \right\},
  \qquad
  \widehat M
  =
  \mathbb{P}_n
  \left\{
    h(V)\,b(V)\,b(V)^\top
  \right\}.
\]

\begin{theorem}[Asymptotic linearity of the working-model projection]
\label{thm:msm-projection-al}
Suppose Assumptions~\ref{assump:consistency}--\ref{assump:baseline-v} hold,
\(M\) is nonsingular, \(h\) and \(b\) are bounded, the pseudo-outcomes are
constructed with cross-fitting, the cumulative estimated ratios are bounded in
probability, and the nuisance product-rate condition
\[
  \sum_{s=0}^{\tau}
  \left\|
    \widehat r_s^{\d}-r_{s}^{\d}
  \right\|_{2}
  \left\|
    \widehat m_s^{\d}-m_{s}^{\d}
  \right\|_{2}
  =
  o_p(n^{-1/2})
\]
holds in the marginal \(L_2(P_0)\) norm. Then
\[
  \widehat\beta_{\d}
  -
  \beta_{\d}
  =
  \mathbb{P}_n D_{\beta,\d}
  +
  o_p(n^{-1/2}),
  \qquad
  \sqrt n
  \left(
    \widehat\beta_{\d}
    -
    \beta_{\d}
  \right)
  \rightsquigarrow
  N(0,\Sigma_\beta),
  \qquad
  \Sigma_\beta
  =
  \E
  \left\{
    D_{\beta,\d}
    D_{\beta,\d}^\top
  \right\},
\]
and \(\Sigma_\beta\) is consistently estimated by the sandwich estimator
\[
  \widehat\Sigma_\beta
  =
  \widehat M^{-1}
  \,
  \mathbb{P}_n
  \left[
    h(V)^2\,
    b(V)\,b(V)^\top
    \left\{
      \widehat\varphi_{\d}(O)
      -
      b(V)^\top
      \widehat\beta_{\d}
    \right\}^2
  \right]
  \widehat M^{-1}.
\]
The same statements hold for the contrast with \(\varphi_\tau\),
\(D_{\beta,\tau}\), and \(\beta_\tau\).
\end{theorem}

\begin{proof}
Writing
\(\widehat\beta_{\d}-\beta_{\d}
=\widehat M^{-1}\mathbb{P}_n\,h\,b\,
\{\widehat\varphi_{\d}-b^\top\beta_{\d}\}\)
and decomposing
\(\widehat\varphi_{\d}
=\varphi_{\d}(O;\eta_0^{\d})
+b_{\widehat\eta}^{\d}(V)+e_{\widehat\eta}(O)\)
as in the proof of Theorem~\ref{thm:dr-learner-oracle}, the term involving
\(\varphi_{\d}(O;\eta_0^{\d})\) yields the influence-function expansion, the
conditionally mean-zero term \(e_{\widehat\eta}\) is
\(o_p(n^{-1/2})\) by cross-fitting, and the bias term is bounded by
\(\|h\,b\|_\infty\,\E|b_{\widehat\eta}^{\d}(V)|\), which is
\(o_p(n^{-1/2})\) under the product-rate condition by
Proposition~\ref{prop:second-order-remainder} and the Cauchy--Schwarz
inequality; \(\widehat M\to M\) in probability completes the argument. Details
are given in the Supplementary Material.
\end{proof}

Note that the rate condition in Theorem~\ref{thm:msm-projection-al} involves
the \emph{marginal} \(L_2(P_0)\) norms rather than the subgroup-restricted
norms of Corollary~\ref{cor:product-rates}: because the projection averages
over \(V\), no subgroup-specific nuisance accuracy is required.

\begin{remark}[Simultaneous bands for the fitted summary]
\label{rem:msm-simultaneous}
Wald intervals for individual coefficients (for example, a linear-trend slope)
follow directly. For the fitted curve, a simultaneous band over any finite grid
\(v_1,\ldots,v_G\) is
\(b(v_g)^\top\widehat\beta_{\d}\pm
c_{1-\alpha}\{b(v_g)^\top\widehat\Sigma_\beta\,b(v_g)/n\}^{1/2}\), where
\(c_{1-\alpha}\) is the max-\(|Z|\) critical value computed from the Gaussian
vector with correlation matrix induced by
\(\{b(v_g)^\top\widehat\Sigma_\beta\,b(v_{g'})\}_{g,g'}\), exactly as in
Section~\ref{subsec:subgroup-tmle}. Homogeneity of the effect in \(V\) within
the working model corresponds to the null that all non-intercept coordinates of
\(\beta_\tau\) vanish, testable by a Wald statistic based on
\(\widehat\Sigma_\beta\).
\end{remark}

\begin{remark}[The subgroup TMLE as a saturated projection]
\label{rem:msm-saturated}
When \(V\) is discrete with support \(\{v_1,\ldots,v_K\}\), taking
\(h\equiv 1\) and the saturated basis
\(b(v)=\{\mathds{1}(v=v_1),\ldots,\mathds{1}(v=v_K)\}^\top\) gives
\(\widehat\beta_{\d,k}
=\mathbb{P}_n\{\mathds{1}(V=v_k)\widehat\varphi_{\d}\}/\widehat p_{v_k}\), the
one-step subgroup estimator, with \(D_{\beta,\d}\) reducing coordinatewise to
the subgroup efficient influence functions
\(D_{\d,v}\) of Section~\ref{subsec:subgroup-tmle}. The working-model
formulation therefore nests the subgroup analysis, and the TMLE of
Section~\ref{subsec:subgroup-tmle} may be viewed as its targeted refinement,
whose plug-in form guarantees that estimates respect the parameter space.
\end{remark}

\subsection{Undersmoothed B-spline sieves}
\label{subsubsec:msm-sieve}

When \(V\) is continuous and the target is \(\theta_{\d}(v)\) itself, the same
weighted least-squares estimator may be used with a basis whose dimension grows
with \(n\). Let \(V\) have compact support
\(\mathcal V\subset\R^{d_V}\), and let
\(b_{k_n}(v)\in\R^{k_n}\) collect tensor-product B-splines of fixed order
\(\ell\) on a quasi-uniform knot sequence, with \(k_n\to\infty\). Write
\(\widehat\theta_{\d,n}(v)=b_{k_n}(v)^\top\widehat\beta_{\d,n}\) for the
resulting series estimator and
\(\sigma_n^2(v)=b_{k_n}(v)^\top\Sigma_{\beta,n}\,b_{k_n}(v)/n\) for its
pointwise variance, with \(\Sigma_{\beta,n}\) as in
Theorem~\ref{thm:msm-projection-al} at dimension \(k_n\). For local bases such
as B-splines, \(\sigma_n(v)\asymp (k_n/n)^{1/2}\), while the sup-norm
approximation bias of the projection satisfies
\(\sup_v|b_{k_n}(v)^\top\beta_{\d,n}-\theta_{\d}(v)|
=O(k_n^{-s/d_V})\)
when \(\theta_{\d}\) is \(s\)-smooth with \(s\le\ell\). \emph{Undersmoothing}
selects \(k_n\) large enough that the bias is negligible relative to the
pointwise standard deviation, so that the confidence statements center on
\(\theta_{\d}\) rather than on its projection.

\begin{assumption}[Sieve regularity]
\label{assump:sieve}
(i) The density of \(V\) is bounded away from zero and infinity on
\(\mathcal V\); (ii) \(\theta_{\d}\) belongs to a H\"older class of smoothness
\(s\le\ell\) on \(\mathcal V\); (iii) the dimension satisfies
\(k_n\log n/n\to 0\) and the undersmoothing condition
\(n\,k_n^{-(2s/d_V+1)}\to 0\); (iv) the pseudo-outcome has uniformly bounded
conditional \((2+\varepsilon)\)-th moments given \(V\), as holds under
condition~3 of Assumption~\ref{assump:subgroup-tmle-regularity}; and (v) the
nuisance product-rate condition of Theorem~\ref{thm:msm-projection-al} holds.
\end{assumption}

Condition (iii) is nonvacuous whenever
\(n^{d_V/(2s+d_V)}\ll k_n\ll n/\log n\): the sieve dimension is taken of larger
order than the mean-squared-error-optimal choice, sacrificing rate optimality
for centering. Under (i), B-spline Gram matrices have eigenvalues bounded away
from zero and infinity and
\(\sup_v\|b_{k_n}(v)\|\asymp k_n^{1/2}\), so the conditions of the least-squares
series literature apply
\cite{belloni2015series,semenova2021debiased}.

\begin{proposition}[Pointwise and uniform sieve inference]
\label{prop:sieve-inference}
Under Assumptions~\ref{assump:consistency}--\ref{assump:baseline-v}
and~\ref{assump:sieve}, with cross-fitted pseudo-outcomes,
\[
  \frac{
    \widehat\theta_{\d,n}(v)
    -
    \theta_{\d}(v)
  }{
    \widehat\sigma_n(v)
  }
  \rightsquigarrow
  N(0,1)
  \quad\text{for each fixed }v,
\]
and the multiplier-bootstrap critical value
\(\widehat c_{1-\alpha}\) for
\(\sup_{v\in\mathcal V}
|\widehat\theta_{\d,n}(v)-\theta_{\d}(v)|/\widehat\sigma_n(v)\)
yields uniformly valid confidence bands,
\[
  \P
  \left\{
    \theta_{\d}(v)
    \in
    \widehat\theta_{\d,n}(v)
    \pm
    \widehat c_{1-\alpha}\,
    \widehat\sigma_n(v)
    \ \text{for all }
    v\in\mathcal V
  \right\}
  \to
  1-\alpha.
\]
The analogous statements hold for the contrast \(\tau(\cdot)\).
\end{proposition}

\begin{proof}
The result follows by verifying the conditions of
\cite{belloni2015series} and \cite{semenova2021debiased} for least-squares
series regression with the estimated outcome \(\widehat\varphi_{\d}\):
the decomposition
\(\widehat\varphi_{\d}
=\varphi_{\d}(O;\eta_0^{\d})+b_{\widehat\eta}^{\d}(V)+e_{\widehat\eta}(O)\)
reduces the problem to the oracle series regression plus a nuisance-bias
contribution whose projection is uniformly
\(o_p\{(k_n/n)^{1/2}\}\) under condition (v), since the product-rate condition
implies \(\|b_{\widehat\eta}^{\d}\|_{2}=o_p(n^{-1/2})=o_p\{(k_n/n)^{1/2}\}\),
and a conditionally centered contribution controlled by cross-fitting. The
undersmoothing condition (iii) renders the approximation bias
\(o_p\{\sigma_n(v)\}\) uniformly. Verification of the technical conditions is
given in the Supplementary Material.
\end{proof}

\begin{remark}[Nuisance rates are again weaker than root-\(n\)]
\label{rem:sieve-rates}
As in Remark~\ref{rem:dr-learner-rates}, requiring the nuisance products to be
\(o_p(n^{-1/2})\) in condition (v) is convenient but conservative: since the
sieve standard deviation is of order \((k_n/n)^{1/2}\gg n^{-1/2}\), it suffices
that the products be \(o_p\{(k_n/n)^{1/2}\}\), a strictly weaker requirement
that grows weaker as the band is made more flexible.
\end{remark}

\begin{remark}[Practical selection of the sieve]
\label{rem:sieve-practice}
In practice we take cubic B-splines (\(\ell=4\)), place knots at quantiles of
\(V\), select a mean-squared-error-optimal dimension
\(\widehat k\) by cross-validation on the pseudo-outcome regression, and then
undersmooth by a fixed multiple, e.g.\ \(k_n=\lceil c\,\widehat k\rceil\) with
\(c\in[1.5,2]\); the resulting bands are wider than the cross-validated fit's
naive bands, which is the price of centering on \(\theta_{\d}\) rather than a
smoothed surrogate. Alternatives that avoid undersmoothing --- explicit bias
correction in the second-stage regression \cite{takatsu2022debiased} --- apply
here as well by substituting our pseudo-outcome, exactly as in the
point-treatment setting \cite{bonvini2023flexibly}.
\end{remark}

\section{Theoretical Properties}
\label{sec:theory}

We now give large-sample properties of the proposed estimators. The finite
subgroup TMLE is a regular asymptotically linear estimator of the vector of
subgroup-specific LMTP effects under standard cross-fitting and second-order
remainder conditions. The DR-Learner is analyzed as a regression with estimated
pseudo-outcomes. Its behavior is compared to the benchmark regression that would
be obtained if the nuisance functions were known.

Throughout this section, let \(P_0\) denote the true observed data law. For a
policy \(\mathbbm d\), let \(\eta_0^{\mathbbm d}\) denote the corresponding true
nuisance functions, and let
\[
  \varphi_{\mathbbm d}(O;\eta_0^{\mathbbm d})
\]
denote the uncentered LMTP influence-function pseudo-outcome defined in
Section~\ref{subsec:uncentered-eif-pseudooutcomes}. For a contrast of two
policies \(\mathbbm d_1\) and \(\mathbbm d_0\), define
\[
  \varphi_{\tau}(O;\eta_0)
  =
  \varphi_{\mathbbm d_1}(O;\eta_0^{\mathbbm d_1})
  -
  \varphi_{\mathbbm d_0}(O;\eta_0^{\mathbbm d_0}).
\]

\subsection{Asymptotic linearity of the subgroup TMLE}
\label{subsec:theory-subgroup-tmle}

Suppose \(V\) is discrete with finite support
\[
  \mathcal V = \{v_1,\ldots,v_K\}.
\]
Let
\[
  \theta_{\mathbbm d}
  =
  \{\theta_{\mathbbm d}(v_1),\ldots,\theta_{\mathbbm d}(v_K)\}^\top
\]
denote the vector of policy-specific subgroup LMTP means. For a contrast of two
policies, write
\[
  \tau
  =
  \{\tau(v_1),\ldots,\tau(v_K)\}^\top.
\]
The efficient influence function for the policy-specific subgroup mean
\(\theta_{\mathbbm d}(v)\) is
\[
  D_{\mathbbm d,v}(O)
  =
  \frac{\mathds{1}(V=v)}{p_v}
  \left\{
    \varphi_{\mathbbm d}(O;\eta_0^{\mathbbm d})
    -
    \theta_{\mathbbm d}(v)
  \right\},
  \qquad
  p_v = P_0(V=v).
\]
For a contrast,
\[
  D_{\tau,v}(O)
  =
  D_{\mathbbm d_1,v}(O)
  -
  D_{\mathbbm d_0,v}(O).
\]
Let
\[
  D_\tau(O)
  =
  \{D_{\tau,v_1}(O),\ldots,D_{\tau,v_K}(O)\}^\top.
\]

\begin{assumption}[Regularity for subgroup TMLE]
\label{assump:subgroup-tmle-regularity}
The following conditions hold.
\begin{enumerate}
  \item The identification assumptions in
  Section~\ref{subsec:identification-assumptions} hold for each policy under
  consideration.

  \item For each \(v\in\mathcal V\), \(p_v>0\), and
  \(\min_{v\in\mathcal V}p_v\) is bounded away from zero.

  \item The outcome and sequential regressions are uniformly bounded, and the
  cumulative density ratios
  \(\omega_s^{\mathbbm d}(O)=\prod_{t=0}^s r_t^{\mathbbm d}(A_t,H_t)\) are
  uniformly bounded or have sufficiently many bounded moments.

  \item The nuisance functions are estimated using cross-fitting, so that
  nuisance estimation and evaluation of the estimating equations are performed
  on independent samples.

  \item The targeting step solves the empirical efficient influence-function
  equations up to negligible error:
  \[
    \mathbb{P}_n
    \widehat D_{\mathbbm d,v}
    =
    o_p(n^{-1/2})
    \qquad
    \text{for all }v\in\mathcal V.
  \]

  \item The second-order remainder for each policy satisfies
  \[
    R_{\mathbbm d,n}(v)
    =
    o_p(n^{-1/2})
    \qquad
    \text{for all }v\in\mathcal V.
  \]
\end{enumerate}
\end{assumption}

\begin{theorem}[Asymptotic linearity of the subgroup TMLE]
\label{thm:subgroup-tmle-al}
Under Assumption~\ref{assump:subgroup-tmle-regularity},
\[
  \widehat\theta_{\mathbbm d,\mathrm{tmle}}(v)
  -
  \theta_{\mathbbm d}(v)
  =
  \mathbb{P}_n D_{\mathbbm d,v}
  +
  o_p(n^{-1/2})
\]
for each \(v\in\mathcal V\). Hence
\[
  \sqrt n
  \left\{
    \widehat\theta_{\mathbbm d,\mathrm{tmle}}
    -
    \theta_{\mathbbm d}
  \right\}
  \rightsquigarrow
  N(0,\Sigma_{\mathbbm d}),
\]
where
\[
  \Sigma_{\mathbbm d}
  =
  \mathbb{E}_{P_0}
  \left[
    D_{\mathbbm d}(O)D_{\mathbbm d}(O)^\top
  \right].
\]
For a contrast of two policies,
\[
  \widehat\tau_{\mathrm{tmle}}(v)-\tau(v)
  =
  \mathbb{P}_n D_{\tau,v}
  +
  o_p(n^{-1/2}),
\]
and therefore
\[
  \sqrt n
  \left\{
    \widehat\tau_{\mathrm{tmle}}
    -
    \tau
  \right\}
  \rightsquigarrow
  N(0,\Sigma_\tau),
\]
where
\[
  \Sigma_\tau
  =
  \mathbb{E}_{P_0}
  \left[
    D_\tau(O)D_\tau(O)^\top
  \right].
\]
\end{theorem}

\begin{proof}
We give the argument for a single policy. The contrast result follows by
linearity. By the definition of the subgroup parameter and the efficient
influence function,
\[
  \widehat\theta_{\mathbbm d,\mathrm{tmle}}(v)
  -
  \theta_{\mathbbm d}(v)
  =
  (\mathbb{P}_n-P_0)D_{\mathbbm d,v}
  +
  \mathbb{P}_n
  \widehat D_{\mathbbm d,v}
  +
  R_{\mathbbm d,n}(v)
  +
  o_p(n^{-1/2}).
\]
The empirical estimating equation term
\(\mathbb{P}_n\widehat D_{\mathbbm d,v}\) is \(o_p(n^{-1/2})\) by the targeting
condition. The second-order remainder is \(o_p(n^{-1/2})\) by assumption. The
cross-fitting construction removes the need for Donsker restrictions on the
estimated nuisance functions and allows the leading empirical process term to be
evaluated at the fixed true influence function. Thus
\[
  \widehat\theta_{\mathbbm d,\mathrm{tmle}}(v)
  -
  \theta_{\mathbbm d}(v)
  =
  \mathbb{P}_nD_{\mathbbm d,v}
  +
  o_p(n^{-1/2}).
\]
The multivariate central limit theorem gives the joint normal limit over
\(v\in\mathcal V\).
\end{proof}

\begin{remark}
The empirical covariance estimator
\[
  \widehat\Sigma_{\tau}
  =
  \mathbb{P}_n
  \left[
    \widehat D_{\tau}(O)
    \widehat D_{\tau}(O)^\top
  \right]
\]
is consistent for \(\Sigma_\tau\) under the same conditions, provided the
estimated influence functions converge in \(L_2(P_0)\) to their population
limits.
\end{remark}

\subsection{Oracle equivalence of the DR-Learner}
\label{subsec:theory-dr-learner}

We next study the DR-Learner for the conditional LMTP contrast
\[
  \tau(v)
  =
  \mathbb{E}_{P_0}
  \{\varphi_\tau(O;\eta_0)\mid V=v\}.
\]
Let \(\widehat\varphi_{\tau,i}\) be the cross-fitted estimated pseudo-outcome
defined in Section~\ref{subsec:dr-learner}. Let
\[
  \widehat{\mathbb E}_n[Z\mid V=v]
\]
denote a generic second-stage regression estimator applied to an outcome
\(Z\) and evaluated at \(V=v\). This notation may represent, for example, a
series regression, spline regression, local polynomial regression, random forest,
or any other regression method used in the second stage.

The feasible DR-Learner is
\[
  \widehat\tau_{\mathrm{dr}}(v)
  =
  \widehat{\mathbb E}_n
  \left[
    \widehat\varphi_\tau
    \mid V=v
  \right].
\]
The benchmark regression that would be obtained if the nuisance functions were
known is
\[
  \widetilde\tau(v)
  =
  \widehat{\mathbb E}_n
  \left[
    \varphi_\tau(O;\eta_0)
    \mid V=v
  \right].
\]
This benchmark is not available in practice; it is used only to measure the
cost of estimating the nuisance functions.

Define the conditional pseudo-outcome bias
\[
  b_{\widehat\eta}(v)
  =
  \mathbb{E}_{P_0}
  \left[
    \widehat\varphi_\tau(O)
    -
    \varphi_\tau(O;\eta_0)
    \mid
    \widehat\eta,V=v
  \right].
\]
Equivalently, \(b_{\widehat\eta}(v)\) is the regression function of the
pseudo-outcome estimation error on \(V\), conditional on the nuisance training
sample.

\begin{assumption}[Second-stage stability]
\label{assump:second-stage-stability}
We say that \(\widehat{\mathbb E}_n\) is stable at \(V=v\), with respect to a
distance \(d\), if
\[
  \frac{
    \widehat{\mathbb E}_n
    \left[
      \widehat\varphi_\tau(O)
      \mid V=v
    \right]
    -
    \widehat{\mathbb E}_n
    \left[
      \varphi_\tau(O;\eta_0)
      \mid V=v
    \right]
    -
    \widehat{\mathbb E}_n
    \left[
      b_{\widehat\eta}(V)
      \mid V=v
    \right]
  }{
    \sqrt{
      \mathbb E
      \left[
        \left(
          \widehat{\mathbb E}_n
          \left[
            \varphi_\tau(O;\eta_0)
            \mid V=v
          \right]
          -
          \mathbb E_{P_0}
          \left[
            \varphi_\tau(O;\eta_0)
            \mid V=v
          \right]
        \right)^2
      \right]
    }
  }
  \overset{p}{\longrightarrow}
  0
\]
whenever
\[
  d\{\widehat\varphi_\tau,\varphi_\tau(\eta_0)\}
  \overset{p}{\longrightarrow}
  0.
\]
\end{assumption}

\begin{theorem}[Oracle equivalence of the LMTP DR-Learner]
\label{thm:dr-learner-oracle}
Suppose the identification assumptions hold for both policies
\(\mathbbm d_1\) and \(\mathbbm d_0\). Suppose the nuisance functions are
estimated with cross-fitting and that
Assumption~\ref{assump:second-stage-stability} holds. Then
\[
  \widehat\tau_{\mathrm{dr}}(v)
  -
  \widetilde\tau(v)
  =
  \widehat{\mathbb E}_n
  \left[
    b_{\widehat\eta}(V)
    \mid V=v
  \right]
  +
  o_p\{R_n^\star(v)\}.
\]
Consequently,
\[
  \widehat\tau_{\mathrm{dr}}(v)
  -
  \tau(v)
  =
  \widetilde\tau(v)
  -
  \tau(v)
  +
  \widehat{\mathbb E}_n
  \left[
    b_{\widehat\eta}(V)
    \mid V=v
  \right]
  +
  o_p\{R_n^\star(v)\}.
\]
If
\[
  \widehat{\mathbb E}_n
  \left[
    b_{\widehat\eta}(V)
    \mid V=v
  \right]
  =
  o_p\{R_n^\star(v)\},
\]
then
\[
  \widehat\tau_{\mathrm{dr}}(v)
  -
  \widetilde\tau(v)
  =
  o_p\{R_n^\star(v)\}.
\]
Thus the feasible DR-Learner is oracle equivalent at \(v\).
\end{theorem}

\begin{proof}
Conditional on the nuisance training sample, decompose the estimated
pseudo-outcome as
\[
  \widehat\varphi_\tau(O)
  =
  \varphi_\tau(O;\eta_0)
  +
  b_{\widehat\eta}(V)
  +
  e_{\widehat\eta}(O),
\]
where
\[
  \mathbb{E}_{P_0}
  \{e_{\widehat\eta}(O)\mid \widehat\eta,V\}
  =
  0.
\]
Applying the second-stage regression estimator gives
\[
  \widehat\tau_{\mathrm{dr}}(v)
  =
  \widehat{\mathbb E}_n
  \left[
    \varphi_\tau(O;\eta_0)
    +
    b_{\widehat\eta}(V)
    +
    e_{\widehat\eta}(O)
    \mid V=v
  \right].
\]
By Assumption~\ref{assump:second-stage-stability},
\[
  \widehat\tau_{\mathrm{dr}}(v)
  -
  \widehat{\mathbb E}_n
  \left[
    \varphi_\tau(O;\eta_0)
    \mid V=v
  \right]
  =
  \widehat{\mathbb E}_n
  \left[
    b_{\widehat\eta}(V)
    \mid V=v
  \right]
  +
  o_p\{R_n^\star(v)\}.
\]
Since
\[
  \widetilde\tau(v)
  =
  \widehat{\mathbb E}_n
  \left[
    \varphi_\tau(O;\eta_0)
    \mid V=v
  \right],
\]
the first equation of the theorem follows. Adding and subtracting \(\tau(v)\) gives the second
second equation. The oracle equivalence claim follows if the smoothed conditional bias
is negligible relative to the oracle regression error.
\end{proof}

\begin{corollary}[Consistency of the DR-Learner]
\label{cor:dr-learner-consistency}
Suppose the conditions of Theorem~\ref{thm:dr-learner-oracle} hold. 
If the oracle regression is consistent for the truth, and the 
conditional pseudo-outcome bias 
decays to zero, then the oracle regression will be consistent. That is to 
say,
If \( \widetilde\tau(v)-\tau(v)=o_p(1) \) and
\(
  \widehat{\mathbb E}_n
  \left[
    b_{\widehat\eta}(V)
    \mid V=v
  \right]
  =
  o_p(1),
\) then \( \widehat\tau_{\mathrm{dr}}(v)-\tau(v)=o_p(1) \).
\end{corollary}

\begin{proof}
The result follows immediately from the expansion in
Theorem~\ref{thm:dr-learner-oracle}.
\end{proof}

\subsection{Second-order bias and nuisance-rate conditions}
\label{subsec:second-order-bias}

Both estimators require control of a second-order remainder: condition 6 of
Assumption~\ref{assump:subgroup-tmle-regularity} for the subgroup TMLE, and the
conditional-bias condition
\(\widehat{\mathbb E}_n\{b_{\widehat\eta}(V)\mid V=v\}=o_p\{R_n^\star(v)\}\) of
Theorem~\ref{thm:dr-learner-oracle} for the DR-Learner. In this section we give
these remainders an explicit form and derive sufficient nuisance rate
conditions. The key object is the conditional bias of the uncentered
pseudo-outcome under estimated nuisances,
\[
  b_{\widehat\eta}^{\d}(v)
  =
  \mathbb{E}_{P_0}
  \left\{
    \varphi_{\d}(O;\widehat\eta^{\d})
    -
    \varphi_{\d}(O;\eta_0^{\d})
    \mid V=v
  \right\},
\]
where the expectation is over \(O\) with the estimated nuisances
\(\widehat\eta^{\d}\) held fixed (as guaranteed by cross-fitting). Because
\(\mathbb{E}_{P_0}\{\varphi_{\d}(O;\eta_0^{\d})\mid V=v\}=\theta_{\d}(v)\) by
Theorem~\ref{thm:pseudooutcome-identification}, \(b_{\widehat\eta}^{\d}\) is
exactly the conditional analogue of the scalar second-order remainder, and the
subgroup TMLE remainder satisfies
\(R_{\d,n}(v)=b_{\widehat\eta}^{\d}(v)+o_p(n^{-1/2})\) under the conditions of
Remark~\ref{rem:subgroup-proportion} below.

\begin{proposition}[Product form of the conditional remainder]
\label{prop:second-order-remainder}
Let \(\widehat\omega_{s-1}^{\d}(O)=\prod_{k=0}^{s-1}\widehat
r_k^{\d}(A_k,H_k)\) with the convention \(\widehat\omega_{-1}^{\d}\equiv 1\).
Under Assumptions~\ref{assump:consistency}--\ref{assump:baseline-v}, and
provided the estimated ratios and regressions are uniformly bounded,
\[
  b_{\widehat\eta}^{\d}(v)
  =
  -
  \sum_{s=0}^{\tau}
  \mathbb{E}_{P_0}
  \left[
    \widehat\omega_{s-1}^{\d}(O)
    \left\{
      \widehat r_s^{\d}(A_s,H_s)
      -
      r_{s}^{\d}(A_s,H_s)
    \right\}
    \left\{
      \widehat m_s^{\d}(A_s,H_s)
      -
      m_{s}^{\d}(A_s,H_s)
    \right\}
    \,\middle|\,
    V=v
  \right].
\]
\end{proposition}

\begin{proof}
The argument telescopes the sum defining \(\varphi_{\d}\) one time point at a
time, using at each \(s\) the intervention-density identity
\(\mathbb{E}\{f(A_s^{\d},H_s)\mid H_s\}
=\mathbb{E}\{r_s^{\d}(A_s,H_s)f(A_s,H_s)\mid H_s\}\), which holds for any
bounded \(f\) by Assumption~\ref{assump:policy-regularity} and the definition
of \(r_s^{\d}\). Note that the cumulative weight multiplying the \(s\)-th
product is the \emph{estimated} weight \(\widehat\omega_{s-1}^{\d}\), not its
population counterpart. Full details are given in the Supplementary Material.
\end{proof}

Two features of Proposition~\ref{prop:second-order-remainder} drive the rate
conditions used throughout Section~\ref{sec:theory}. First, the remainder is a
sum of \(\tau+1\) terms, each a product of the estimation errors of one density
ratio and its same-time-point sequential regression; there are no cross-time
products. Second, each product enters only through its conditional mean given
\(V=v\), so all rates may be taken with respect to the subgroup-restricted norm
\[
  \|f\|_{2,v}^2
  =
  \mathbb{E}_{P_0}
  \left\{
    f(O)^2 \mid V=v
  \right\}.
\]

\begin{corollary}[Sufficient product-rate conditions]
\label{cor:product-rates}
Suppose \(\max_{s}\sup_o|\widehat\omega_{s-1}^{\d}(o)|=O_p(1)\). Then, by the
Cauchy--Schwarz inequality,
\[
  \left|b_{\widehat\eta}^{\d}(v)\right|
  \;\lesssim\;
  \sum_{s=0}^{\tau}
  \left\|
    \widehat r_s^{\d}-r_{s}^{\d}
  \right\|_{2,v}
  \left\|
    \widehat m_s^{\d}-m_{s}^{\d}
  \right\|_{2,v}.
\]
In particular, if for every \(s=0,\ldots,\tau\),
\(\|\widehat r_s^{\d}-r_{s}^{\d}\|_{2,v}
  \|\widehat m_s^{\d}-m_{s}^{\d}\|_{2,v}
  =o_p(n^{-1/2})\)
--- as holds when each nuisance converges at the \(o_p(n^{-1/4})\) rate
attainable by many machine-learning estimators under smoothness or sparsity ---
then \(R_{\d,n}(v)=o_p(n^{-1/2})\), and condition 6 of
Assumption~\ref{assump:subgroup-tmle-regularity} is satisfied.
\end{corollary}

\begin{remark}[Multiple robustness]
\label{rem:multiple-robustness}
The product form implies \(2^{\tau+1}\)-fold robustness: the conditional bias
at \(v\) vanishes provided that, for \emph{each} time point \(s\) separately,
at least one of \(\widehat r_s^{\d}\) and \(\widehat m_s^{\d}\) is consistent
in \(\|\cdot\|_{2,v}\), with the other merely bounded. Consistent estimation of
all ratios, or of all sequential regressions, are two special cases.
\end{remark}

\begin{remark}[Weaker conditions for the DR-Learner]
\label{rem:dr-learner-rates}
For the DR-Learner, the requirement of Theorem~\ref{thm:dr-learner-oracle} is
that the smoothed conditional bias be \(o_p\{R_n^\star(v)\}\), where
\(R_n^\star(v)\) is the oracle rate of the second-stage regression. When \(V\)
is continuous, \(R_n^\star(v)\) is slower than \(n^{-1/2}\), so the nuisance
product rates need only be faster than the oracle rate rather than faster than
\(n^{-1/2}\). Oracle equivalence for conditional LMTP effects is therefore
available under strictly weaker nuisance conditions than root-\(n\) inference
for the subgroup vector, in direct analogy with the point-treatment DR-Learner
\cite{kennedy2023towards}.
\end{remark}

\begin{remark}[Subgroup-proportion remainder]
\label{rem:subgroup-proportion}
The subgroup functional \(\theta_{\d}(v)
=\mathbb{E}\{\mathds{1}(V=v)\,m_0^{\d}(A_0^{\d},L_0)\}/p_v\) contributes an
additional remainder through \(\widehat p_v\). A direct expansion of the ratio
shows this contribution is
\(-(\widehat p_v-p_v)\{\widehat\theta_{\d}(v)-\theta_{\d}(v)\}/p_v
+o_p(n^{-1/2})\), a product of a root-\(n\) term and a consistent-estimation
error. It is therefore \(o_p(n^{-1/2})\) whenever
\(\widehat\theta_{\d}(v)\) is consistent and \(p_v\) is bounded away from zero
(condition 2 of Assumption~\ref{assump:subgroup-tmle-regularity}), and is
dominated by the terms in Proposition~\ref{prop:second-order-remainder}.
\end{remark}


\section{Simulation Study}
\label{sec:simulation-study}

We evaluate the subgroup TMLE and its simultaneous-inference procedures in a
three-timepoint LMTP setting (\(\tau=2\)) in which the truth, the sequential
regressions, and the density ratios are all available in closed form, and we
evaluate the DR-Learner in a companion design with a continuous baseline
modifier. Closed-form nuisances serve two purposes: they permit an oracle arm
that anchors every comparison at the efficiency bound, and they make the truth
exact rather than approximated by Monte Carlo.

\subsection{Design: subgroup inference with three exposure timepoints}
\label{subsec:sim1-design}

Baseline covariates are \(L_0=(X,W_1,W_2)\) with \(X\sim N(0,I_{d})\) and
\(W_1,W_2\) independent Bernoulli\((1/2)\) effect modifiers. Time-varying
covariates \(Z_1,Z_2\) and exposures \(A_0,A_1,A_2\) follow structural
equations that are \emph{linear in all endogenous variables} with constant
coefficients, while arbitrary (mean-centered) nonlinearity in \(X\) enters
additively:
\begin{align*}
  A_0 &= \mu(X) + \varepsilon_0,
  &
  Z_1 &= 0.4\,A_0 + q(X) + \nu_1,
  \\
  A_1 &= 0.3\,A_0 + 0.3\,Z_1 + \mu(X) + \varepsilon_1,
  &
  Z_2 &= 0.4\,A_1 + 0.2\,Z_1 + q(X) + \nu_2,
  \\
  A_2 &= 0.3\,A_1 + 0.3\,Z_2 + \mu(X) + \varepsilon_2,
  &
  Y &= f(X) + \beta(X,W)\textstyle\sum_{t=0}^{2}A_t + 0.3\,(Z_1+Z_2)
      + \varepsilon_Y,
\end{align*}
with all noise terms independent standard normal and
\(\beta(X,W)=1+b_1W_1+b_2W_2\) (plus \(0.5\sin(\pi X_1)\) in the complex
setting below). The policy is the additive shift
\(d_t(a_t,h_t)=a_t+\delta\) with \(\delta=1/2\) at every timepoint, and the
target is the vector of shifted means \(\theta_{\d}(v)\) over the four
\emph{overlapping} subgroups
\(\mathcal V=\{W_1=1,\;W_2=1,\;W_1=W_2=1,\;\text{overall}\}\). Overlap is
deliberate: it induces substantial off-diagonal correlation in
\(\Sigma_{\d}\), which is precisely the structure that separates the
max-\(|Z|\) band from Bonferroni.

Because the endogenous coefficients are constants, the shift propagates
deterministically through the system:
\(\Delta_{A_0}=0.5\), \(\Delta_{Z_1}=0.2\), \(\Delta_{A_1}=0.71\),
\(\Delta_{Z_2}=0.324\), \(\Delta_{A_2}=0.8102\), so that, exactly,
\[
  \theta_{\d}(v) - \mathbb{E}(Y\mid V=v)
  =
  \mathbb{E}\{\beta(X,W)\mid V=v\}\times 2.0202
  \;+\;
  0.3\times 1.048/2
  =
  2.0202\,\mathbb{E}\{\beta\mid V=v\} + 0.1572,
\]
and the level terms vanish by mean-centering, so in the linear settings
\(\theta_{\d}(v)\) equals this expression exactly. The closed-form sequential
regressions \(m_t^{\d}\) and Gaussian-shift ratios
\(r_t^{\d}(a_t,h_t)=\exp\{\delta(a_t-\mu_t(h_t))-\delta^2/2\}\) supply the
oracle nuisance arm and the per-subgroup efficiency bounds, which are computed
by evaluating the efficient influence function on a single large Monte Carlo
sample (\(2\times 10^6\) draws) and verified against the exact truths.

We cross four data-generating settings with an estimator roster. The settings
vary the modification strength \((b_1,b_2)\in\{(0,0),(0.15,-0.15),(0.5,-0.5)\}\)
with \(d=2\) and linear \(\mu,q,f\) (so that suitably \(W\)-interacted
main-terms generalized linear models are \emph{exactly} correct), plus a
complex setting with \(d=4\), nonlinear \(\mu,q,f\), and
\(\beta\) depending on \(X\) (so those models are misspecified by
construction). The roster varies: (i) estimator --- one-step versus TMLE, the
latter with the weighted or linear fluctuation of
Section~\ref{subsec:subgroup-tmle}; (ii) nuisance estimation --- correctly
specified or misspecified parametric regressions, an oracle arm evaluating the
closed-form nuisances, and a super learner arm in which each nuisance is fit
once on the full sample and predictions are taken out-of-fold, with candidate
learners refit without each observation's fold and combined using pooled
ensemble weights, the cross-validated targeting construction of
\cite{zheng2011cvtmle}; and (iii) density-ratio estimation --- direct
conditional-density modeling versus the odds-ratio classification
representation. Each cell uses 1{,}000 replications at
\(n\in\{1000,2000,4000,8000\}\)
(\(n\le 4000\) for super-learner arms). Endpoints are \(n\times\)MSE relative
to the efficiency bound, \(\sqrt n\,\)bias, standard-error calibration,
marginal coverage, familywise coverage of the three simultaneous procedures
(unadjusted, Bonferroni, and max-\(|Z|\) based on \(\widehat\Sigma_{\d}\)),
size and power of the homogeneity Wald test with full versus diagonal
covariance, and coverage and power for the nested contrast
\(\theta_{\d}(W_1{=}1)-\theta_{\d}(W_1{=}W_2{=}1)\) using joint versus
independence-based standard errors.

\subsection{Results}
\label{subsec:sim1-results}

Figure~\ref{fig:sim1-simultaneous} displays the headline endpoint. The oracle
arms attain the efficiency bound with nominal marginal and familywise coverage
at all sample sizes, validating the implementation. Unadjusted per-subgroup
intervals undercover as a joint statement in every setting (familywise
coverage near {\color{Mahogany}0.8XX}), while Bonferroni is conservative. The
max-\(|Z|\) band achieves nominal familywise coverage using a critical value of
{\color{Mahogany}X.XX} against Bonferroni's \(2.50\)
({\color{Mahogany}X.X}\% narrower intervals) at a mean off-diagonal influence
function correlation of {\color{Mahogany}0.XX}. In the linear settings the
correctly specified parametric arms are efficient, and the classification-based
ratio arm is numerically indistinguishable from the density-based arm; in the
complex setting the misspecified parametric arms exhibit
\(\sqrt n\)-growing bias and degrading coverage while the super-learner arms
remain near the bound, with coverage {\color{Mahogany}0.9XX} at
\(n=4000\). The homogeneity test with the full covariance holds its size on the
null setting ({\color{Mahogany}0.0XX} at nominal 0.05) and dominates the
diagonal-covariance version in power, and the joint-covariance standard error
for the nested contrast is correctly calibrated where the independence-based
standard error is conservative by a factor of {\color{Mahogany}X.X}, forfeiting
corresponding power. Complete results for all settings, arms, and endpoints
appear in the Supplementary Material.

\begin{figure}[H]
\centering
% \includegraphics[width=0.95\textwidth]{figures/figM1_manuscript_simultaneous_coverage.pdf}
\caption{Familywise coverage of nominal 95\% simultaneous intervals for the
four overlapping subgroup LMTP means (\(\tau=2\), additive shift
\(\delta=1/2\)), for the targeted estimator with super-learner nuisances, under
a setting with no effect modification (left) and the complex nonlinear setting
with strong modification (right). The max-\(|Z|\) band calibrated to the
estimated joint influence-function covariance attains nominal familywise
coverage with narrower intervals than Bonferroni; the oracle-nuisance
max-\(|Z|\) curve is shown as a reference. Error bars are Monte Carlo 95\%
intervals over 1{,}000 replications.}
\label{fig:sim1-simultaneous}
\end{figure}

\subsection{Design and results: DR-Learner with a continuous modifier}
\label{subsec:sim2}

The second study retains the longitudinal structure above but replaces the
binary modifiers with a continuous baseline variable \(V\sim\mathrm{Unif}(-1,1)\),
correlated with the confounders through
{\color{Mahogany}[finalize: \(V=\Phi^{-1}\)-coupled to \(X_1\) at correlation
0.3]}, and sets \(\beta(V)=1+\gamma\,\{\sin(\pi V)+V\}\), a smooth nonmonotone
modification profile, so that
\(\theta_{\d}(v)=2.0202\,\beta(v)+0.1572\) exactly. We compare the DR-Learner
of Section~\ref{subsec:dr-learner} --- cross-fitted pseudo-outcomes
\(\widehat\varphi_{\d}\) regressed on \(V\) by local linear smoothing with
cross-validated bandwidth --- against (i) the oracle second-stage regression
that uses \(\varphi_{\d}(O;\eta_0)\), and (ii) a plug-in learner that regresses
\(\widehat m_0^{\d}(A_0^{\d},L_0)\) on \(V\) without the influence-function
correction. Endpoints are integrated root mean squared error over the central
90\% of the support of \(V\) and pointwise coverage of 95\% confidence bands.
Consistent with Theorem~\ref{thm:dr-learner-oracle}, the DR-Learner's RMSE
tracks the oracle within {\color{Mahogany}X}\% by \(n=4000\) even with
super-learner nuisances, while the plug-in learner inherits first-order
nuisance bias in the complex setting
(Supplementary Figure~{\color{Mahogany}SX}).

\section{Application}
\label{sec:application}

We apply the proposed methods to estimate the effect of deforestation on
malaria incidence in Madagascar over 2010--2022, using monthly confirmed case
counts reported by 2{,}773 public health clinics linked to satellite-derived
annual forest-cover loss within each clinic's catchment area
\cite{vieilledent2018combining}. Madagascar spans strongly contrasting
ecosystems --- humid forests, dry deciduous forests, spiny thickets, succulent
woodlands, and subhumid forests --- and the ecological dynamics linking land
use change to vector habitat differ across them, making effect heterogeneity
by ecological zone a primary scientific question rather than a secondary
subgroup analysis. Each clinic is assigned to one of \(K=5\) zones
\cite{dinerstein2017ecoregion}, with the sparsely populated mangrove and
montane regions folded into their surrounding zones.

The observed data for clinic \(i\) are
\(O_i=(L_{0,i},A_{0,i},L_{1,i},A_{1,i},\ldots,A_{\tau,i},Y_i)\) with annual
exposure timepoints (\(\tau=12\)): \(A_{t,i}\) is forest loss in catchment
\(i\) in year \(2010+t\), \(L_{t,i}\) collects time-varying confounders ---
annual precipitation, minimum, maximum, and range of temperature, soil
moisture, the oceanic Ni\~no index, population size and growth, crop
productivity, and lagged malaria incidence --- and \(Y_i\) is the clinic-level
malaria case rate in 2022. Baseline covariates \(L_{0,i}\) additionally include
elevation, ecological zone \(V_i\), and small-area poverty estimates. We
consider the policy \(\d\) that reduces observed deforestation by
{\color{Mahogany}10}\% at every timepoint,
\(d_t(a_t,h_t)=(1-{\color{Mahogany}0.10})\,a_t\), an intervention that respects
the observed exposure distribution in each zone rather than assigning a common
fixed exposure, and we report the contrast against the natural course. We
assume interference across catchments is negligible after conditioning on each
clinic's own history; analyses accommodating spillover between neighboring
health sheds are the subject of ongoing work and are discussed in
Section~\ref{sec:discussion}.

All sequential regressions and density ratios were estimated with a super
learner over {\color{Mahogany}[list: penalized GLMs with Poisson loss,
multivariate adaptive regression splines, random forests, gradient-boosted
trees]}, using out-of-fold predictions, with count outcomes weighted by
catchment person-time. The joint subgroup TMLE with the weighted fluctuation
targeted the six-dimensional parameter (five zone-specific means and the
national mean) for each policy. As a positivity diagnostic we tracked the
effective sample size implied by the cumulative estimated ratios
\(\widehat\omega_{s}^{\d}\) at every timepoint, which fell from \(n=2{,}773\)
to {\color{Mahogany}X{,}XXX} at \(s=\tau\); ratios were truncated at
{\color{Mahogany}[bound]} affecting {\color{Mahogany}X.X}\% of clinic-years.

Table~{\color{Mahogany}\ref{tab:application-results}} reports the estimated
policy effects with 95\% pointwise and max-\(|Z|\) simultaneous intervals. The
estimated national effect of the {\color{Mahogany}10}\% reduction policy is a
change of {\color{Mahogany}XX.X} (95\% CI: {\color{Mahogany}XX.X, XX.X}) cases
per 1{,}000 person-years. Zone-specific estimates ranged from
{\color{Mahogany}XX.X} in the {\color{Mahogany}[zone]} to
{\color{Mahogany}XX.X} in the {\color{Mahogany}[zone]}; the simultaneous
max-\(|Z|\) critical value was {\color{Mahogany}X.XX}, against a Bonferroni
value of \(2.64\) for six parameters, with mean off-diagonal influence-function
correlation {\color{Mahogany}0.XX} induced by the nesting of the national
parameter. The homogeneity Wald test with the full estimated covariance
{\color{Mahogany}[rejected / did not reject]} equality of zone-specific effects
(\(\chi^2_{4}={\color{Mahogany}XX.X}\),
\(p={\color{Mahogany}0.XXX}\)). To characterize heterogeneity in continuous
baseline features, we applied the DR-Learner with \(V\) taken to be baseline
forest cover, regressing the pseudo-outcome contrast on \(V\) by local linear
smoothing; the estimated conditional effect curve
(Figure~{\color{Mahogany}\ref{fig:application-drlearner}}) indicates
{\color{Mahogany}[shape: e.g., effects concentrated among catchments with high
baseline cover, plateauing above the XXth percentile]}, consistent with the
frontier-clearing mechanisms hypothesized in the ecological literature
\cite{macdonald2019amazon}. Sensitivity analyses using an alternative
forest-loss product and alternative malaria endpoints, together with all
diagnostics, appear in the Supplementary Material.

\section{Discussion}
\label{sec:discussion}

We have developed efficient, multiply robust methodology for effect
heterogeneity analysis of longitudinal modified treatment policies: a jointly
targeted estimator of subgroup-specific LMTP means and contrasts with
simultaneous inference calibrated to the estimated joint influence-function
covariance, and a DR-Learner for conditional LMTP effect functions of
continuous baseline covariates with oracle-equivalence guarantees. The two
procedures are complementary answers to the same scientific question at
different resolutions of the modifier, and both reduce to nuisance objects ---
sequential regressions and policy density ratios --- that are already computed
by standard LMTP software, so heterogeneity analysis imposes essentially no
additional nuisance-estimation burden beyond the marginal analysis.

Three aspects of the joint approach merit emphasis. First, targeting the
subgroup vector jointly, rather than refitting within subgroups, stabilizes
nuisance estimation for small subgroups and yields the full covariance matrix
\(\widehat\Sigma\) as a byproduct; the simultaneous max-\(|Z|\) band and the
homogeneity test are then available at no extra cost. The width advantage of
the max-\(|Z|\) band over Bonferroni grows with the correlation among subgroup
estimators, which is substantial precisely in the designs applied researchers
favor --- overlapping subgroups, nested contrasts against an overall effect,
and marginal-plus-subgroup reporting as in our application. Second, the
product form of the second-order remainder
(Proposition~\ref{prop:second-order-remainder}) shows that the sequential
double robustness familiar from marginal LMTP estimation
\cite{diaz2023lmtp} carries over conditionally on \(V\), time point by time
point. Third, for continuous modifiers the DR-Learner requires only that
nuisance error products beat the oracle nonparametric rate, not root-\(n\), so
flexible machine learning remains admissible exactly where it is most needed.

Several limitations point to future work. Our analysis treats clinic
catchments as non-interfering units after conditioning on their own histories;
deforestation and malaria plausibly spill over between neighboring health
sheds, and extending conditional LMTP estimands and estimators to accommodate
network or spatial interference is a natural next step. Second, the DR-Learner
outputs a nonparametric curve; interpretation tools developed for
point-treatment conditional effects --- additive projections, partial
dependence summaries, and variable-importance measures for effect modification
\cite{bonvini2023flexibly,semenova2021debiased,hines2022variable} --- extend
directly by substituting our pseudo-outcome, and projections of
\(\theta_{\d}(\cdot)\) onto low-dimensional working models summaries 
with root-\(n\) inference. Third, subgroup-specific
\emph{policies}, in which different shifts \(\d_v\) are assigned to different
subgroups --- of interest, for example, when feasible reductions in
deforestation differ by ecological zone, or when immune-marker shifts
attainable by a refined vaccine differ across subpopulations
\cite{hejazi2023stochastic} --- are immediately accommodated by our framework,
since the subgroup TMLE targets each \(\theta_{\d_v}(v)\) jointly with any
collection of policies. Finally, count outcomes recorded against person-time
invite loss functions and fluctuation submodels adapted to rate parameters;
our weighted fluctuation is a first step, and a systematic treatment of
nonparametric rate targets under LMTPs would be valuable.

\section{References}
\bibliography{refs}

\section{Acknowledgments}

\end{document}
